% ------------------------------------------------------------------------
%  My Format for a Article
% ------------------------------------------------------------------------
\documentclass[12pt]{article}
\oddsidemargin 0.0in
\textwidth 6.5in
\textheight 9in
\topmargin -0.35in
\headheight 0.0in
\hfuzz 10pt

\renewcommand{\thefootnote}{\fnsymbol{footnote}}
\newcommand{\be}{\begin{equation}}
\newcommand{\ee}{\end{equation}}
\newcommand{\Curl}{\nabla \times}
\newcommand{\Div}{\nabla \cdot}

% ------------------------------------------------------------------------
% Definitions so that LaTeX won't throw a
% "Too many unprocessed floats" tantrum
% ------------------------------------------------------------------------
\renewcommand{\topfraction}{.85}
\renewcommand{\bottomfraction}{.7}
\renewcommand{\textfraction}{.15}
\renewcommand{\floatpagefraction}{.66}
\renewcommand{\dbltopfraction}{.66}
\renewcommand{\dblfloatpagefraction}{.66}
\setcounter{topnumber}{9}
\setcounter{bottomnumber}{9}
\setcounter{totalnumber}{20}
\setcounter{dbltopnumber}{9}

% ------------------------------------------------------------------------
% Packages to use for a dvips document:
% ------------------------------------------------------------------------
\usepackage{pslatex}
\usepackage[]{graphicx}
\usepackage[]{hyperref}
\begin{document}

% ------------------------------------------------------------------------
% Title Page:
% ------------------------------------------------------------------------
\thispagestyle{empty}

\vspace*{2.5in}

\vspace*{0.5in}
\begin{center}
{\LARGE Finite Element Method Magnetics: OctaveFEMM}


\vspace*{16pt} {\Large User's Manual}

\vspace*{16pt} \today

\vspace*{48pt} David Meeker \\
\href{mailto:dmeeker@ieee.org}{\tt dmeeker@ieee.org} \\
\href{http://www.femm.info}{\tt http://www.femm.info} \\
\end{center}

\newpage


% ------------------------------------------------------------------------
% Main Text:
% ------------------------------------------------------------------------
\section{Introduction}

OctaveFEMM is a Matlab toolbox that allows for the operation of Finite
Element Method Magnetics (FEMM) via a set of Matlab functions.  The toolbox
works with Octave, a Matlab clone.

When OctaveFEMM starts up a FEMM process, the usual FEMM user interface
is displayed and is fully functional.  The user then has the choice of
accomplishing modeling and analysis tasks either exclusively through
functions implemented by the toolbox, or by a combination of manual and
programmatic operations -- whichever is easiest for the task at hand.

The syntax of the OctaveFEMM toolbox closely mirrors that of FEMM's existing
Lua scripting language interface associated with FEMM v4.2.  However, there
are some differences between the Lua functions and the analogous Octave/Matlab
implementations:
\begin{itemize}
\item All strings are enclosed in single quotes, rather than
double quotes as in Lua.
\item Functions in Lua that have no arguments require a set of
empty parentheses after the function name ({\em e.g.} {\tt
mi\_analyze()}).  In Octave or Matlab, no parentheses should be used needed ({\em e.g.} {\tt
mi\_analyze} with the OctaveFEMM toolbox).
\item Several commands have also been added to OctaveFEMM that have no
analog in Lua.  These commands streamline the drawing of new
geometries with the OctaveFEMM toolbox, as well is the collection
of data from solutions.
\end{itemize}

Perhaps the most remarkable difference between Lua and OctaveFEMM,
however, is due to the matrix-oriented nature of Octave/Matlab.  In just about any
OctaveFEMM function in which it would be desiable to enter an array
of points such that multiple copies of an operation are performed, OctaveFEMM
will correctly interpret the input perform the requested operation
on every element in the array.  In addition, for any function in which
the coordinates of a point are required, that point can be
specified as an array with two elements instead of specifying each
element separately.  In functions that require the specification
of multiple points, those points can be entered as an array of
two-element arrays.

\section{Installation and Startup}

\subsection{Installation for Matlab and Octave 3}
The OctaveFEMM distribution is automatically installed with FEMM 4.2 in the
\verb+mfiles+ subdirectory.  The absolute directory is typically
\verb+c:\Program Files\femm42\mfiles+.  To use OctaveFEMM with Octave or
Matlab, this path needs to be added to the program's search path. To add
this path to the search path, type the following lines at the Matlab or
Octave 3.X.X command prompt:

\begin{verbatim}
	addpath('c:\\femm42\\mfiles');
	savepath;
\end{verbatim}

Alternatively, in Matlab, the interactive {\tt pathtool} command can be used
to add the {\tt mfiles} directory to the search path.

\subsection{Startup}
To start an OctaveFEMM session, use the {\tt openfemm} function.
This function starts up a FEMM process to which OctaveFEMM calls
are sent. When done with OctaveFEMM, the FEMM process can be shut
down via the {\tt closefemm} function.

A number of examples that use OctaveFEMM to analyze various problems
are included in the directory \verb+cd c:\Program Files\femm42\examples+

%-----------------------------------------------------------------------------------
%
%  OctaveFEMM Documentation
%
%-----------------------------------------------------------------------------------

% Magnetics-specific functions
% ------------------------------------------------------------------------
%  My Format for a Article

\section{Common Command Set}

There are a number of FEMM-specific Octave that are not associated with any particular problem type.
These functions manipulate the appearance of the main window and
other top-level components like the Lua console and Point
Properties output window.

\begin{itemize}

\item {\tt openfemm(bHide)} The {\tt bHide} parameter determines whether or not the main FEMM window is visible when FEMM is started. In many cases, it is useful to use {\tt openfemm(1)} to hide the main window while a script is running.  In hidden window mode, analysis calls like \verb+mi_analyze+ hide the analysis window by default.

\item {\tt clearconsole} Clears the output window of the Lua console.

\item {\tt newdocument(doctype)} Creates a new preprocessor document
and opens up a new preprocessor window.  Specify {\tt doctype} to
be {\tt 0} for a magnetics problem, {\tt 1} for an electrostatics
problem, {\tt 2} for a heat flow problem, or {\tt 3} for a current
flow problem. Alternative syntax for this function is {\tt create(doctype)}

\item{\tt hideconsole} Hides the floating Lua console window.

\item{\tt hidepointprops} Hides the floating FEMM Properties display window.

\item \texttt{messagebox('message')} displays the \texttt{'message' }string to the
screen in a pop-up message box.

\item {\tt opendocument('filename')} Opens a document specified by {\tt filename}.

\item \texttt{print(item1,item2,...)} This is standard Lua ``print'' command directed to the
output of the Lua console window. Any number of comma-separated items can be
printed at once via the print command.

\item \texttt{prompt('message')} This function allows a FEMM script to prompt a
user for input. When this command is used, a dialog box pops up with the
\texttt{'message'} string on the title bar of the dialog box. The user can
enter in a single line of input via the dialog box.  Prompt returns the
user's input to Octave and parses it using Octave's {\tt eval}
command.

\item{\tt showconsole} Displays the floating Lua console window.

\item{\tt showpointprops} Displays the floating FEMM Properties
display window.

\item{\tt smartmesh(state)} Calling with a state of 0 turns off ``smart mesh'' functionality for the present session; calling with a state of 1 turns ``smarth meshing'' on. The {\tt smartmesh} function applies across all problem types. 
To set smart meshing on or off on a file-by-file basis, use  \verb+mi_smartmesh+,\verb+ei_smartmesh+, \verb+hi_smartmesh+, and \verb+ci_smartmesh+, depending on problem type.  Note that calling the global \verb+smartmesh+ with a value of -1 sets the program to defer to the file-by-file setting rather than forcing all smart meshing on or off.

\item{\tt main\_minimize} minimizes the main FEMM window.

\item{\tt main\_maximize} maximizes the main FEMM window.

\item{\tt main\_restore} restores the main FEMM window from a
 minimized or maximized state.

\item{\tt main\_resize(width,height)} resizes the main FEMM
 window client area to width $\times$ height.

\end{itemize}


\section{Magnetics Preprocessor Command Set}

A number of different commands are available in the preprocessor.

\subsection{Object Add/Remove Commands}

\begin{itemize}

\item{\tt mi\_addnode(x,y)} Add a new node at x,y

\item{\tt mi\_addsegment(x1,y1,x2,y2)} Add a new line segment from node
closest to (x1,y1) to node closest to (x2,y2)

\item{\tt mi\_addblocklabel(x,y)} Add a new block label at (x,y)

\item{\tt mi\_addarc(x1,y1,x2,y2,angle,maxseg)} Add a new arc segment
from the nearest node to (x1,y1) to the nearest node to (x2,y2)
with angle `angle' divided into `maxseg' segments.

\item{\tt mi\_drawline(x1,y1,x2,y2)} Adds nodes at (x1,y1) and
(x2,y2) and adds a line between the nodes.

\item{\tt mi\_drawpolyline([x1,y1;x2,y2'...])} Adds nodes at each
of the specified points and connects them with segments.

\item{\tt mi\_drawpolygon([x1,y1;x2,y2'...])} Adds nodes at each
of the specified points and connects them with segments to form a
closed contour.

\item{\tt mi\_drawarc(x1,y1,x2,y2,angle,maxseg)} Adds nodes at
(x1,y1) and (x2,y2) and adds an arc of the specified angle and
discretization connecting the nodes.

\item{\tt mi\_drawrectangle(x1,y1,x2,y2)} Adds nodes at the
corners of a rectangle defined by the points (x1,y1) and (x2,y2),
then adds segments connecting the corners of the rectangle.

\item{\tt mi\_deleteselected} Delete all selected objects.

\item{\tt mi\_deleteselectednodes} Delete selected nodes.

\item{\tt mi\_deleteselectedlabels} Delete selected block labels.

\item{\tt mi\_deleteselectedsegments} Delete selected segments.

\item{\tt mi\_deleteselectedarcsegments} Delete selects arcs.
\end{itemize}

\subsection{Geometry Selection Commands}

\begin{itemize}
\item{\tt mi\_clearselected} Clear all selected nodes, blocks, segments
and arc segments.

\item{\tt mi\_selectsegment(x,y)} Select the line segment closest to
(x,y)

\item{\tt mi\_selectnode(x,y)} Select the node closest to (x,y).
Returns the coordinates of the selected node.

\item{\tt mi\_selectlabel(x,y)} Select the label closet to (x,y).
Returns the coordinates of the selected label.

\item{\tt mi\_selectarcsegment(x,y)} Select the arc segment closest to
(x,y)

\item{\tt mi\_selectgroup(n)} Select the $n^{th}$ group of nodes, segments, arc
segments and blocklabels. This function will clear all previously selected
elements and leave the editmode in 4 (group)

\item{\tt mi\_selectcircle(x,y,R,editmode)} selects objects within a circle of radius
R centered at (x,y).  If only x, y, and R paramters are given, the current
edit mode is used.  If the editmode parameter is used, 0 denotes nodes, 1 denotes segments,
2 denotes block labels, 3 denotes arcs, and 4 specifies that all entity types are to be selected.

\item{\tt mi\_selectrectangle(x1,y1,x2,y2,editmode)} selects objects within a rectangle
defined by points (x1,y1) and (x2,y2). If no editmode parameter is supplied,
the current edit mode is used.  If the editmode parameter is used, 0 denotes nodes, 1 denotes segments,
2 denotes block labels, 3 denotes arcs, and 4 specifies that all entity types are to be selected.

\end{itemize}

\subsection{Object Labeling Commands}

\begin{itemize}
\item{\tt mi\_setnodeprop('propname',groupno)} Set the selected nodes to
have the nodal property {\tt 'propname'} and group number {\tt groupno}.

\item{\tt mi\_setblockprop('blockname', automesh, meshsize, 'incircuit',
magdir, group, turns)} Set the selected block labels to have the
properties:
\begin{itemize}
\item Block property {\tt 'blockname'}.
\item {\tt automesh}: 0 = mesher defers to mesh size constraint defined in {\tt meshsize},
         1 = mesher automatically chooses the mesh density.
\item {\tt meshsize}: size constraint on the mesh in the block marked by this label.
\item Block is a member of the circuit named {\tt 'incircuit'}
\item The magnetization is directed along an angle in measured in degrees denoted by the parameter
        {\tt magdir}
\item A member of group number {\tt group}
\item The number of turns associated with this label is denoted by {\tt turns}.
\end{itemize}

\item{\tt mi\_setsegmentprop('propname', elementsize, automesh, hide,
group)} Set the selected segments to have:
\begin{itemize}
\item Boundary property {\tt 'propname'}
\item Local element size along segment no greater than {\tt
elementsize}
\item {\tt automesh}:  0 = mesher defers to the element constraint defined by {\tt elementsize},
        1 = mesher automatically chooses mesh size along the selected segments
\item {\tt hide}: 0 =  not hidden in post-processor, 1 == hidden in post processor
\item A member of group number {\tt group}
\end{itemize}

\item{\tt mi\_setarcsegmentprop(maxsegdeg, 'propname', hide, group)} Set the
selected arc segments to:
\begin{itemize}
\item Meshed with elements that span at most {\tt maxsegdeg} degrees per
element
\item Boundary property {\tt 'propname'}
\item {\tt hide}: 0 =  not hidden in post-processor, 1 == hidden in post processor
\item A member of group number {\tt group}
\end{itemize}

\item{\tt mi\_setgroup(n)} Set the group associated of the selected items to n.

\end{itemize}

\subsection{Problem Commands}

\begin{itemize}
\item{\tt mi\_probdef(freq,units,type,precision,depth,minangle,(acsolver))} changes the
problem definition. Set {\tt freq} to the desired frequency in
Hertz.  The {\tt units} parameter specifies the units used for
measuring length in the problem domain.  Valid {\tt 'units'}
entries are {\tt 'inches'}, {\tt 'millimeters'}, {\tt
'centimeters'}, {\tt 'mils'}, {\tt 'meters'}, and {\tt
'micrometers'}. Set the parameter {\tt problemtype} to {\tt 'planar'} for a 2-D
planar problem, or to {\tt 'axi'} for an axisymmetric problem. The
{\tt precision} parameter dictates the precision required by the
solver.  For example, entering {\tt 1E-8} requires the RMS of the
residual to be less than $10^{-8}$.  A fifth parameter, representing the depth of
the problem in the into-the-page direction for 2-D planar problems.  Specify the depth
to be zero for axisymmetric problems. The sixth parameter represents the minimum angle
constraint sent to the mesh generator -- 30 degrees is the usual
choice for this parameter.  The acsolver parameter specifies which solver is to be
used for AC problems: 0 for successive approximation, 1 for Newton.

\item{\tt mi\_analyze(flag)}
runs the magnetics solver.  The {\tt flag} parameter
controls whether the solver window is visible or minimized.  For a
visible window,  specify {\tt 0}. For a minimized window, {\tt flag} should be set to {\tt 1}.
If no value is specified for {\tt flag}, the visibility of the solver is inherited from the main window,
{\em i.e.} if the main window is minimized, the solver runs minimized, too.

\item{\tt mi\_loadsolution} loads and displays the solution corresponding to the
current geometry.

\item {\tt mi\_setfocus('documentname')} Switches the
magnetics input file upon which commands are to act. If
more than one magnetics input file is being edited at a time,
this command can be used to switch between files so that the
mutiple files can be operated upon programmatically. {\tt
'documentname'} should contain the name of the desired document as
it appears on the window's title bar.

\item{\tt mi\_saveas('filename')} saves the file with name {\tt 'filename'}.

\end{itemize}

\subsection{Mesh Commands}

\begin{itemize}
\item{\tt mi\_createmesh} runs the mesher to create a mesh. Note that this is not a
necessary precursor of performing an analysis, as {\tt
mi\_analyze} will make sure the mesh is up to date before running
an analysis. The number of elements in the mesh is pushed back onto
the lua stack.

\item{\tt mi\_showmesh} shows the mesh.

\item{\tt mi\_purgemesh} clears the mesh out of both the screen and
memory.
\end{itemize}

\subsection{Editing Commands}

\begin{itemize}
\item{\tt mi\_copyrotate(bx, by, angle, copies )}
        \begin{itemize}
        \item{\tt bx, by} -- base point for rotation
        \item{\tt angle} -- angle by which the selected objects are incrementally
        shifted to make each copy.  {\tt angle} is measured in degrees.
        \item{\tt copies} -- number of copies to be produced from
        the selected objects.
        \end{itemize}

\item{\tt mi\_copyrotate2(bx, by, angle, copies, editaction )}
        \begin{itemize}
        \item{\tt bx, by} -- base point for rotation
        \item{\tt angle} -- angle by which the selected objects are incrementally
        shifted to make each copy.  {\tt angle} is measured in degrees.
        \item{\tt copies} -- number of copies to be produced from
        the selected objects.
        \item{\tt editaction}  0 --nodes, 1 -- lines (segments), 2 --block labels, 3 -- arc
                segments, 4- group
        \end{itemize}

\item{\tt mi\_copytranslate(dx, dy, copies)}
        \begin{itemize}
        \item{\tt dx,dy} -- distance by which the selected objects are incrementally shifted.
        \item{\tt copies} -- number of copies to be produced from the selected objects.
        \end{itemize}

\item{\tt mi\_copytranslate2(dx, dy, copies, editaction)}
        \begin{itemize}
        \item{\tt dx,dy} -- distance by which the selected objects are incrementally shifted.
        \item{\tt copies} -- number of copies to be produced from the selected objects.
        \item{\tt editaction}  0 --nodes, 1 -- lines (segments), 2 --block labels, 3 -- arc
                segments, 4- group
        \end{itemize}

\item{\tt mi\_createradius(x, y, r)} turns a corner located at {\tt (x,y)} into a curve of radius {\tt r}.

\item{\tt mi\_moverotate(bx,by,shiftangle)}
        \begin{itemize}
        \item{\tt bx, by} -- base point for rotation
        \item{\tt shiftangle} -- angle in degrees by which the selected objects are rotated.
        \end{itemize}

\item{\tt mi\_moverotate2(bx,by,shiftangle, editaction)}
        \begin{itemize}
        \item{\tt bx, by} -- base point for rotation
        \item{\tt shiftangle} -- angle in degrees by which the selected objects are rotated.
        \item{\tt editaction}  0 --nodes, 1 -- lines (segments), 2 --block labels, 3 -- arc
                segments, 4- group
        \end{itemize}

\item{\tt mi\_movetranslate(dx,dy)}
        \begin{itemize}
        \item{\tt dx,dy} -- distance by which the selected objects are shifted.
        \end{itemize}

\item{\tt mi\_movetranslate2(dx,dy,editaction)}
        \begin{itemize}
        \item{\tt dx,dy} -- distance by which the selected objects are shifted.
        \item{\tt editaction}  0 --nodes, 1 -- lines (segments), 2 --block labels, 3 -- arc
                segments, 4- group
        \end{itemize}

\item{\tt mi\_scale(bx,by,scalefactor)}
        \begin{itemize}
        \item{\tt bx, by} -- base point for scaling
        \item{\tt scalefactor} -- a multiplier that determines how
        much the selected objects are scaled
        \end{itemize}

\item{\tt mi\_scale2(bx,by,scalefactor,editaction)}
        \begin{itemize}
        \item{\tt bx, by} -- base point for scaling
        \item{\tt scalefactor} -- a multiplier that determines how
        much the selected objects are scaled
        \item{\tt editaction}  0 --nodes, 1 -- lines (segments), 2 --block labels, 3 -- arc
                segments, 4- group
        \end{itemize}

\item{\tt mi\_mirror(x1,y1,x2,y2)}
mirror the selected objects about a line passing through the points
{\tt (x1,y1)} and {\tt (x2,y2)}.

\item{\tt mi\_mirror2(x1,y1,x2,y2,editaction)}
mirror the selected objects about a line passing through the points
{\tt (x1,y1)} and {\tt (x2,y2)}. Valid {\tt editaction} entries are
0 for nodes, 1 for lines (segments), 2 for block labels, 3 for arc
segments, and 4 for groups.


\item{\tt mi\_seteditmode(editmode)}
Sets the current editmode to:
        \begin{itemize}
        \item{\tt 'nodes'} - nodes
        \item{\tt 'segments'} - line segments
        \item{\tt 'arcsegments'} - arc segments
        \item{\tt 'blocks'} - block labels
        \item{\tt 'group'} - selected group
        \end{itemize}
This command will affect all subsequent uses of the other editing
commands, if they are used WITHOUT the {\tt editaction} parameter.
\end{itemize}

\subsection{Zoom Commands}

\begin{itemize}
\item{\tt mi\_zoomnatural} zooms to a ``natural'' view with sensible extents.
\item{\tt mi\_zoomout} zooms out by a factor of 50\%.
\item{\tt mi\_zoomin} zoom in by a factor of 200\%.
\item{\tt mi\_zoom(x1,y1,x2,y2)}
Set the display area to be from the bottom left corner specified by
{\tt (x1,y1}) to the top right corner specified by {\tt (x2,y2)}.
\end{itemize}

\subsection{View Commands}

\begin{itemize}
\item{\verb+mi_showgrid+} Show the grid points.
\item{\verb+mi_hidegrid+} Hide the grid points points.
\item{\verb+mi_gridsnap('flag')+}
Setting {\tt flag} to 'on' turns on snap to grid, setting {\tt
flag} to {\tt 'off'} turns off snap to grid.
\item{\tt mi\_snapgridon} turns on snap-to-grid. Equivalent to {\tt mi\_gridsnap('on')}.

\item{\tt mi\_snapgridoff} turns off snap-to-grid. Equivalent to {\tt mi\_gridsnap('off')}.

\item{\verb+mi_setgrid(density,'type')+} Change the grid spacing.  The {\tt density}
parameter specifies the space between grid points, and the {\tt
type} parameter is set to {\tt 'cart'} for cartesian coordinates or
{\tt 'polar'} for polar coordinates.
\item{\tt mi\_refreshview} Redraws the current view.
\item{\tt mi\_minimize} minimizes the active magnetics input view.
\item{\tt mi\_maximize} maximizes the active magnetics input view.
\item{\tt mi\_restore} restores the active magnetics input view from a
 minimized or maximized state.
\item{\tt mi\_resize(width,height)} resizes the active magnetics input
 window client area to width $\times$ height.
\item{\tt mi\_getview} grabs the currently displayed magnetics input view and
 returns it as an image matrix, in the form \texttt{imread} would return.
\end{itemize}

\subsection{Object Properties}

\begin{itemize}
\item \texttt{mi\_getmaterial('materialname')} fetches the material specified by \texttt{materialname} 
from the materials library.

\item{\tt mi\_addmaterial('matname', mu{\_}x, mu{\_}y, H{\_}c,
J, Cduct, Lam{\_}d, Phi{\_}hmax, lam{\_}fill, LamType, Phi{\_}hx, Phi{\_}hy, nstr, dwire)} adds a
new material with called {\tt 'matname'} with the material
properties:
        \begin{itemize}
        \item{\tt mu{\_}x} Relative permeability in the x- or r-direction.
        \item{\tt mu{\_}y} Relative permeability in the y- or z-direction.
        \item{\tt H{\_}c} Permanent magnet coercivity in
        Amps/Meter.
        \item{\tt J} Applied source current
        density in Amps/mm$^2$.
        \item{\tt Cduct} Electrical conductivity of the material
        in~MS/m.
        \item{\tt Lam{\_}d} Lamination thickness in millimeters.
        \item{\tt Phi\_hmax} Hysteresis lag angle in degrees, used for nonlinear BH curves.
        \item{\tt Lam{\_}fill} Fraction of the volume occupied per lamination that
        is actually filled with iron (Note that this parameter defaults to 1 in the
        {\tt femm} preprocessor dialog box because, by default, iron completely
        fills the volume)
        \item{\tt Lamtype} Set to
                \begin{itemize}
                \item 0 -- Not laminated or laminated in plane
                \item 1 -- laminated x or r
                \item 2 -- laminated y or z
                \item 3 -- magnet wire
                \item 4 -- plain stranded wire
                \item 5 -- Litz wire
                \item 6 -- square wire
                \item 7 -- 10\% copper-clad aluminum (CCA) wire
                \item 8 -- 15\% copper-clad aluminum (CCA) wire
                \end{itemize}
        \item{\tt Phi\_hx} Hysteresis lag in degrees in the x-direction for linear problems.
        \item{\tt Phi\_hy} Hysteresis lag in degrees in the y-direction for linear problems.
		\item{\tt nstr} Number of strands in the wire build. Should be 1 for Magnet or Square wire.
		\item{\tt dwire} Diameter of each of the wire's constituent strand in millimeters.
        \end{itemize}
Note that not all properties need be defined -- properties that aren't defined are assigned default
values.


\item{\tt mi\_addbhpoint('blockname',b,h)} Adds a B-H data point the
the material specified by the string {\tt 'blockname'}.  The point to be added
has a flux density of {\tt b} in units of Teslas and a field
intensity of {\tt h} in units of Amps/Meter.

\item{\tt mi\_addbhpoints('blockname',points)} adds a whole B-H curve to the material named {\tt 'blockname'}, where {\tt points} is an $n$-by-2 matrix whose columns are $b$ in Teslas and $h$ in Amps/Meter. The three-argument form {\tt mi\_addbhpoints('blockname',b,h)} takes the two columns separately. Calls {\tt mi\_addbhpoint} once per row.

\item{\tt mi\_clearbhpoints('blockname')} Clears all B-H data points
associatied with the material specified by {\tt 'blockname'}.

\item{\tt mi\_addpointprop('pointpropname',a,j)}
adds a new point property of name {\tt 'pointpropname'} with either
a specified potential {\tt a} in units Webers/Meter
or a point current {\tt j} in units of Amps. Set the
unused parameter pairs to 0.

\item{\tt mi\_addboundprop('propname', A0, A1, A2, Phi, Mu, Sig, c0, c1,
BdryFormat, ia, oa)} \\ adds a new boundary property with name {\tt
'propname'}
        \begin{itemize}
        \item For a ``Prescribed A'' type boundary condition, set the {\tt A0,
        A1, A2} and {\tt Phi} parameters as required. Set all other
        parameters to zero.

        \item For a ``Small Skin Depth'' type boundary condtion, set the {\tt Mu}
        to the desired relative permeability and {\tt Sig} to the desired
        conductivity in MS/m.  Set {\tt BdryFormat} to 1 and all other
        parameters to zero.

        \item To obtain a ``Mixed'' type boundary condition, set {\tt C1} and
        {\tt C0} as required and {\tt BdryFormat} to 2.  Set all other
        parameters to zero.

        \item For a ``Strategic dual image'' boundary, set {\tt BdryFormat} to 3
        and set all other parameters to zero.

        \item For a ``Periodic'' boundary condition, set {\tt BdryFormat} to 4 and
        set all other parameters to zero.

        \item For an ``Anti-Perodic'' boundary condition, set {\tt BdryFormat} to
        5 set all other parameters to zero.
        
        \item  For a ``Periodic Air Gap'', set BdryFormat to 6. Parameters {\tt ia} and {\tt oa} specify the inner and outer  boundary angles, respectively.

        \item  For an ``Anti-periodic Air Gap'', set BdryFormat to 7.  The same  {\tt ia} and {\tt oa} parameters also apply here."

        \end{itemize}

\item{\tt mi\_addcircprop('circuitname', i, circuittype)} \\
adds a new circuit property with name {\tt 'circuitname'} with a prescribed current.
The {\tt circuittype} parameter is 0 for a parallel-connected circuit and 1 for a
series-connected circuit.
\item{\tt mi\_deletematerial('materialname')} deletes the material named {\tt
'materialname'}.
\item{\tt mi\_deleteboundprop('propname')} deletes the boundary property named
{\tt 'propname'}.
\item{\tt mi\_deletecircuit('circuitname')} deletes the circuit named {\tt
circuitname}.
\item{\tt mi\_deletepointprop('pointpropname')} deletes the point property named
{\tt 'pointpropname'}
\item{\verb+mi_modifymaterial('BlockName',propnum,value)+} This
function allows for modification of a material's properties without
redefining the entire material ({\em e.g.} so that current can be
modified from run to run).  The material to be modified is
specified by {\tt 'BlockName'}.  The next parameter is the number
of the property to be set. The last number is the value to be
applied to the specified property.  The various properties that can
be modified are listed below:
\begin{center}
\begin{tabular}{lll} \hline
{\tt propnum}& Symbol & Description \\ \hline
 0 & {\tt BlockName} & Name of the material \\
 1 & $\mu_x$ & x (or r) direction relative permeability \\
 2 & $\mu_y$ & y (or z) direction relative permeability \\
 3 & $H_c$   & Coercivity, Amps/Meter \\
 4 & $J$   & Source current density, MA/m$^2$ \\
 5 & $\sigma$ & Electrical conductivity, MS/m \\
 6 & $d_{lam}$  & Lamination thickness, mm \\
 7 & $\phi_{hmax}$ & Hysteresis lag angle for nonlinear problems, degrees \\
 8 & LamFill & Iron fill fraction \\
 9 & LamType & 0 = None/In plane, 1 = parallel to x, 2=parallel to y \\
 10 & $\phi_{hx}$ & Hysteresis lag in x-direction for linear problems, degrees \\
 11 & $\phi_{hy}$ & Hysteresis lag in y-direction for linear problems, degrees \\
 \hline
 \end{tabular}
 \end{center}
\item{\verb+mi_modifyboundprop('BdryName',propnum,value)+}
This function allows for modification of a boundary property. The
BC to be modified is specified by {\tt 'BdryName'}.  The next
parameter is the number of the property to be set. The last number
is the value to be applied to the specified property.  The various
properties that can be modified are listed below:
\begin{center}
\begin{tabular}{lll} \hline
{\tt propnum}& Symbol & Description \\ \hline
 0 & {\tt BdryName} & Name of boundary property \\
 1 & $A_0$ & Prescribed A parameter \\
 2 & $A_1$ & Prescribed A parameter \\
 3 & $A_2$ & Prescribed A parameter \\
 4 & $\phi$ & Prescribed A phase \\
 5 & $\mu$ & Small skin depth relative permeability \\
 6 & $\sigma$ & Small skin depth conductivity, MS/m \\
 7 & $c_0$ & Mixed BC parameter \\
 8 & $c_1$ & Mixed BC parameter \\
 9 & {\tt BdryFormat} & Type of boundary condition: \\
   &                 & 0 = Prescribed A \\
   &                 & 1 = Small skin depth \\
   &                 & 2 = Mixed \\
   &                 & 3 = Strategic Dual Image \\
   &                 & 4 = Periodic \\
   &                 & 5 = Antiperiodic \\ 
   &                 & 6 = Periodic Air Gap \\ 
   &                 & 7 = Antiperiodic Air Gap \\ 
10 & $ia$ &  Inner boundary angle for air gap element \\
11 & $oa$ & Outer boundary angle for air gap element \\ \hline
\end{tabular}
\end{center}
\item{\verb+mi_modifypointprop('PointName',propnum,value)+}
This function allows for modification of a point property. The
point property to be modified is specified by {\tt 'PointName'}.
The next parameter is the number of the property to be set. The
last number is the value to be applied to the specified property.
The various properties that can be modified are listed below:
\begin{center}
\begin{tabular}{lll} \hline
{\tt propnum}& Symbol & Description \\ \hline
 0 & {\tt PointName} & Name of the point property \\
 1 & $A$ & Nodal potential, Weber/Meter \\
 2 & $J$ & Nodal current, Amps\\
 \hline
\end{tabular}
\end{center}
\item{\verb+mi_modifycircprop('CircName',propnum,value)+}
This function allows for modification of a circuit property. The
circuit property to be modified is specified by {\tt 'CircName'}.
The next parameter is the number of the property to be set. The
last number is the value to be applied to the specified property.
The various properties that can be modified are listed below:
\begin{center}
\begin{tabular}{lll} \hline
{\tt propnum}& Symbol & Description \\ \hline
 0 & {\tt CircName} & Name of the circuit property \\
 1 & $i$ & Total current \\
 2 & {\tt CircType} & 0 = Parallel, 1 = Series \\ \hline
 \end{tabular}
 \end{center}
\item{\tt mi\_setcurrent('CircName',i)} sets the current in the
circuit specified by {\tt 'CircName'} to {\tt i}.
\end{itemize}

\subsection{Miscellaneous}
\begin{itemize}
\item{\tt mi\_savebitmap('filename')} saves a bitmapped screenshot of the current
view to the file specified by {\tt 'filename'}.
\item{\tt mi\_savemetafile('filename')} saves a metafile screenshot of the current
view to the file specified by {\tt 'filename'}.
\item{\tt mi\_refreshview} Redraws the current view.
\item{\tt mi\_close} Closes current magnetics preprocessor
document and destroys magnetics preprocessor window.
\item{\tt mi\_shownames(flag)} This function allow the user to display or hide the block label
names on screen.  To hide the block label names, {\tt flag} should be 0.  To display the
names, the parameter should be set to 1.
\item{\tt mi\_hidenames} hides the block label names. Equivalent to {\tt mi\_shownames(0)}.

\item{\tt mi\_gettitle} returns the file path of the current document as a string.

\item{\tt mi\_setcomment('comment')} sets the comment/description string associated with the current problem.

\item{\tt mi\_readdxf('filename')} This function imports a dxf file specified by {\tt 'filename'}.
\item{\tt mi\_savedxf('filename')} This function saves geometry information in a dxf file specified by {\tt 'filename'}.
\item{\tt mi\_defineouterspace(Zo,Ro,Ri)} defines
an axisymmetric external region to be used in conjuction with the
Kelvin Transformation method of modeling unbounded problems.  The
{\tt Zo} parameter is the z-location of the origin of the outer region,
the {\tt Ro} parameter is the radius of the outer region, and the {\tt
Ri} parameter is the radius of the inner region ({\em i.e.} the region of
interest). In the exterior region, the permeability varies as a function of
distance from the origin of the external region.  These parameters
are necessary to define the permeability variation in the external
region.
\item{\tt mi\_attachouterspace} marks all
selected block labels as members of the external region used for
modeling unbounded axisymmetric problems via the Kelvin
Transformation.
\item{\tt mi\_detachouterspace} undefines all selected block labels
as members of the external region used for modeling unbounded axisymmetric
problems via the Kelvin Transformation.

\item{\tt mi\_attachdefault()} marks the
selected block label as the default block label.  This block label
is applied to any region that has not been explicitly labeled.

\item{\tt mi\_detachdefault()} undefines the default
attribute for the selected block labels.

\item{\tt mi\_makeABC(n,R,x,y,bc)} creates a series of circular shells that emulate the
impedance of an unbounded domain (i.e. an Improvised Asymptotic Boundary
Condition).  The {\tt n} parameter contains the number of shells to be used
(should be between 1 and 10), {\tt R} is the radius of the solution domain, and
{\tt (x,y)} denotes the center of the solution domain.  The {\tt bc} parameter should
be specified as 0 for a Dirichlet outer edge or 1 for a Neumann outer edge.
If the function is called without all the parameters, the function makes up
reasonable values for the missing parameters.

\item{\tt mi\_setprevious(filename,prevtype)} defines the previous solution to be used as the basis for an AC incremental
permeability or frozen permeability solution.  The {\tt prevtype} field is an integer that specifies whether the solution 
is to be incremental permeability ({\tt 1}) or frozen permeability ({\tt 2}). The filename should include the {\tt .ans}
extension, {\em e.g.} {\tt mi\_setprevious('mymodel.ans',1)}

\end{itemize}

\section{Magnetics Post Processor Command Set}

There are a number of scripting commands designed to operate in
the postprocessing environment.

\subsection{Data Extraction Commands}

\begin{itemize}
\item{\tt mo\_lineintegral(type)} Calculate the line integral for the defined contour
\begin{small}\begin{center}
\begin{tabular}{llllll}\hline
 {\tt type} & name & values 1 & values 2 & values 3 & values 4 \\  \hline
 0 & B.n & total B.n & avg B.n & - & -\\
 1 & H.t & total H.t & avg H.t & - & - \\
 2 & Contour length & surface area & - & -\\
 3 & Stress Tensor Force & DC r/x force & DC y/z force & $2\times$ r/x force & $2\times$ y/z force \\
 4 & Stress Tensor Torque& DC torque & $2\times$ torque & - & - \\
 5 & (B.n)\^{}2 & total (B.n)\^{}2 & avg (B.n)\^{}2 & - & - \\ \hline
\end{tabular}
\end{center}
\end{small}
Returns typically two values. The first
value is the result of the integral calculation, and
the second value is the average of the quantity of interest over the contour.
The only exception is integral 3, which evaluates Maxwell's stress
tensor.  This integral can return up to four results. For force
and torque results, the $2\times$ results are only relevant for
problems where $\omega \neq 0$.

\item{\tt mo\_blockintegral(type)}
Calculate a block integral for the selected blocks
\begin{center}
\begin{tabular}{ll} \hline
 Type & Definition \\ \hline
 0 & $A \cdot J$ \\
 1 & A \\
 2 & Magnetic field energy \\
 3 & Hysteresis and/or lamination losses \\
 4 & Resistive losses \\
 5 & Block cross-section area \\
 6 & Total losses \\
 7 & Total current \\
 8 & Integral of $B_x$ (or $B_r$) over block \\
 9 & Integral of $B_y$ (or $B_z$) over block \\
 10 & Block volume \\
 11 & x (or r) part of steady-state Lorentz force \\
 12 & y (or z) part of steady-state Lorentz force \\
 13 & x (or r) part of $2\times$ Lorentz force \\
 14 & y (or z) part of $2\times$ Lorentz force \\
 15 & Steady-state Lorentz torque \\
 16 & $2 \times$ component of Lorentz torque \\
 17 & Magnetic field coenergy \\
 18 & x (or r) part of steady-state weighted stress tensor force \\
 19 & y (or z) part of steady-state weighted stress tensor force \\
 20 & x (or r) part of $2\times$ weighted stress tensor force \\
 21 & y (or z) part of $2\times$ weighted stress tensor force \\
 22 & Steady-state weighted stress tensor torque \\
 23 & $2 \times$ component of weighted stress tensor torque \\
 24 & $R^2$ ({\em i.e.} moment of inertia / density) \\ 
 25 & x (or r) part of $1\times$ weighted stress tensor force \\
 26 & y (or z) part of $1\times$ weighted stress tensor force \\
 27 & $1 \times$ component of weighted stress tensor torque \\
 28 & x (or r) part of $1\times$ Lorentz force \\
 29 & y (or z) part of $1\times$ Lorentz force \\
 30 & $1 \times$ component of Lorentz torque \\ \hline
 \end{tabular}
 \end{center}

\item{\tt mo\_getpointvalues(x,y)}
Get the values associated with the point at (x,y).  The function returns an array whose contents are, in order:
\begin{center}
\begin{tabular}{ll} \hline
Symbol & Definition \\ \hline
 A & Potential A or flux $\phi$ \\
 B1 &  $B_x$ if planar, $B_r$ if axisymmetric \\
 B2  & $B_y$ if planar, $B_z$ if axisymmetric \\
 Sig & conductivity $\sigma$ \\
 E & stored energy density\\
 H1 & $H_x$ if planar, $H_r$ if axisymmetric \\
 H2 &  $H_y$ if planar, $H_z$ if axisymmetric \\
 Je & eddy current density \\
 Js & source current density\\
 Mu1 & $\mu_x$ if planar, $\mu_r$ if axisymmetric \\
 Mu2 & $\mu_y$ if planar, $\mu_z$ if axisymmetric \\
 Pe & Power density dissipated through ohmic losses \\
 Ph & Power density dissipated by hysteresis \\
 ff & Fill factor: lamination fill for laminated materials, winding fill for wire materials, otherwise $-1$ (solid region) or 1 (multi-turn region) \\ \hline
 \end{tabular}
\end{center}
The following series of functions retrieves smaller subsets of
these results.

\item{\tt mo\_geta(x,y)} Get the potential associated with the point at (x,y).
For planar problems, the reported potential is vector potential
$A$.  For axisymmetric problems, $2 \pi r A$ is reported.

\item{\tt mo\_getb(x,y)} Get the magnetic flux density associated with the point at (x,y).
The return value is a list with two element representing $B_x$ and
$B_y$ for planar problems and $B_r$ and $B_z$ for axisymmetric
problems.

\item{\tt mo\_getconductivity(x,y)} Gets the conductivity associated with the point at (x,y).

\item{\tt mo\_getenergydensity(x,y)} Gets the magnetic field energy density associated with the point at (x,y).

\item{\tt mo\_geth(x,y)} Get the magnetic field intensity associated with the point at (x,y).
The return value is a list with two element representing $H_x$ and
$H_y$ for planar problems and $H_r$ and $H_z$ for axisymmetric
problems.

\item{\tt mo\_getj(x,y)} Get the electric current density associated with the point at (x,y).

\item{\tt mo\_getmu(x,y)} Get the relative magnetic permeability associated with the point at (x,y).
The return value is a list with two element representing $\mu_x$ and
$\mu_y$ for planar problems and $\mu_r$ and $\mu_z$ for axisymmetric
problems.

\item{\tt mo\_getpe(x,y)} Get the ohmic loss density associated with the point at (x,y).

\item{\tt mo\_getph(x,y)} Get the hysteresis/laminated eddy current loss density associated with the point at (x,y).

\item{\tt mo\_getfill(x,y)} Get the fill factor associated with the point at (x,y): for a wire material, the winding fill ({\em i.e.} the average fraction of the volume filled with conductor); for a laminated material, the lamination fill; otherwise $-1$ for a solid region or 1 for a multi-turn region.

\item \verb+mo_makeplot(PlotType,NumPoints,Filename,FileFormat)+
Allows Octave access to FEMM's X-Y plot routines.  If only {\tt PlotType} or only {\tt PlotType}
and {\tt NumPoints} are specified, the command is interpreted as a request to plot the
requested plot type to the screen.  If, in addition, the {\tt Filename} parameter is specified,
the plot is instead written to disk to the specified file name as an extended metafile.
If the {\tt FileFormat} parameter is also, the command is instead interpreted as a command to
write the data to disk to the specfied file name, rather than display it to make a
graphical plot.
Valid entries for {\tt PlotType} are:
\begin{center}
\begin{tabular}{ll} \hline
{\tt PlotType} & Definition \\ \hline
0 & Potential \\
1 & $|B|$ \\
2 & $B \cdot n$ \\
3 & $B \cdot t$ \\
4 & $|H|$ \\
5 & $H \cdot n$ \\
6 & $H \cdot t$ \\
7 & $J_{eddy}$ \\
8 & $J_{source}+J_{eddy}$ \\
\hline
\end{tabular}
\end{center}
Valid file formats are
\begin{center}
\begin{tabular}{ll} \hline
{\tt FileFormat} & Definition \\ \hline
0 & Multi-column text with legend \\
1 & Multi-column text with no legend \\
2 & Mathematica-style formatting \\
\hline
\end{tabular}
\end{center}
For example, if one wanted to plot $B \cdot n$ to the screen with 200 points evaluated to
make the graph, the command would be:

\begin{tabular}{l} {\tt mo\_makeplot(2,200)} \end{tabular}

If this plot were to be written to disk as a metafile, the command would be:

\begin{tabular}{l} \verb+mo_makeplot(2,200,'c:\temp\myfile.emf')+ \end{tabular}

To write data instead of a plot to disk, the command would
be of the form:

\begin{tabular}{l} \verb+mo_makeplot(2,200,'c:\temp\myfile.txt',0)+ \end{tabular}

\item \verb+mo_getprobleminfo+
Returns info on problem description.  Returns four values:
\begin{center}
\begin{tabular}{ll} \hline
Return value & Definition \\ \hline 
1 &  problem type  \\ 
2 & frequency in Hz  \\
3 & problem depth in meters \\ 
4 & length unit used to draw the problem, represented in meters\\ \hline
\end{tabular}
\end{center}

\item{\verb+mo_getcircuitproperties('circuit')+}
Used primarily to obtain impedance information associated with
circuit properties.  Properties are returned for the circuit
property named {\tt 'circuit'}. Six values are returned by the
function.  In order, these parameters are:
\begin{itemize}
        \item{\verb+current+} Current carried by the circuit.
        \item{\verb+volts+}  Voltage drop across the circuit in the circuit.
        \item{\verb+flux+} Circuit's flux linkage
\end{itemize}




\item{\verb+mo_getgapb("BdryName",angle)+}
Computes the radial and tangential flux density on the centerline of the
specified air gap element name (\verb+BdryName+) at the specified \verb+angle+.  The angle is specified in degrees.

\item{\verb+mo_getgapa("BdryName",angle)+}
Computes the magnetic vector potential on the centerline of the
specified air gap element name (\verb+BdryName+) at the specified \verb+angle+.  The angle is specified in degrees.

\item{\verb+mo_gapintegral("BdryName",inttype)+}
Computes an integral specifiedy by \verb+inttype+ over an air gap element specified
by \verb+BdryName+. Values for {\tt inttype} in are:
\begin{center}
\begin{tabular}{ll} \hline
0  & DC Torque\\
1  & DC Force\\
2  & Stored Energy\\
3  & 2X Torque\\
4  & 2X Force\\
5  & Interaction Torque\\
6  & Interaction Force\\
\hline
\end{tabular}
\end{center}
Torque and energy integrals return one result; force integrals return two results for the x- and y- direction force, respectively.

\item{\verb+mo_getgapharmonics("BdryName",n)+}
Returns the {\tt acc}, {\tt acs},{\tt brc},{\tt  brs},{\tt  btc}, and {\tt bts}.  These quantities represent the
components of vector potential, radial flux density, and tangential flux density on the centerline of the specified air gap element at the specified angle.  
For the $n^{th}$ harmonic, vector potential, radial flux density, and tangential flux density can be represented explicitly as functions of angle via:
\begin{eqnarray}
A    & = & acc \cos{n \theta} + acs \sin{n \theta} \nonumber  \\ 
B_r & = & brc \cos{n \theta}+ brs \sin{n \theta}   \nonumber \\
B_\theta & = & btc \cos{n \theta} + bts \sin{n \theta}  \nonumber
\end{eqnarray}
If the function is called with just the {\tt BdryName}, the function returns the number of harmonics available.
 
\end{itemize}

\subsection{Selection Commands}
\begin{itemize}
\item{\tt mo\_seteditmode(mode)} Sets the mode of the postprocessor to
point, contour, or area mode.  Valid entries for {\tt mode} are
{\tt 'point'}, {\tt 'contour'}, and {\tt 'area'}.
\item{\tt mo\_selectblock(x,y)} Select the block that contains point (x,y).
\item{\tt mo\_groupselectblock(n)} Selects all of the blocks that are labeled by block
        labels that are members of group {\tt n}. If no number is specified ({\em i.e.} {\tt mo\_groupselectblock} ),
                all blocks are selected.

\item{\tt mo\_addcontour(x,y)} Adds a contour point at (x,y). If this
is the first point then it starts a contour, if there are existing
points the contour runs from the previous point to this point. The
{\tt mo\_addcontour} command has the same functionality as a
right-button-click contour point addition when the program is
running in interactive mode.
\item{\tt mo\_bendcontour(angle,anglestep)} Replaces the straight line
formed by the last two points in the contour by an arc that spans {\tt angle}
degrees.  The arc is actually composed of many straight lines, each
of which is constrained to span no more than {\tt anglestep} degrees.
The {\tt angle} parameter can take on values from -180 to 180 degrees.
The {\tt anglestep} parameter must be greater than zero.  If there are less
than two points defined in the contour, this command is ignored.
\item{\tt mo\_selectpoint(x,y)} Adds a contour point at the closest
input point to (x,y).  If the selected point and a previous
selected points lie at the ends of an arcsegment, a contour is
added that traces along the arcsegment.  The {\tt mo\_selectpoint}
command has the same functionality as the left-button-click contour
point selection when the program is running in interactive mode.
\item{\tt mo\_clearcontour} Clear a prevously defined contour
\item{\tt mo\_clearblock} Clear block selection
\end{itemize}

\subsection{Zoom Commands}
\begin{itemize}
\item{\verb+mo_zoomnatural+} Zoom to the natural boundaries of
the geometry.
\item{\verb+mo_zoomin+} Zoom in one level.
\item{\verb+mo_zoomout+} Zoom out one level.
\item{\tt mo\_zoom(x1,y1,x2,y2)} Zoom to the window defined by lower left corner (x1,y1)
and upper right corner (x2,y2).

\end{itemize}

\subsection{View Commands}
\begin{itemize}
\item{\verb+mo_showmesh+} Show the mesh.
\item{\verb+mo_hidemesh+} Hide the mesh.
\item{\verb+mo_showpoints+} Show the node points from the input geometry.
\item{\verb+mo_hidepoints+} Hide the node points from the input geometry.
\item{\tt mo\_smooth('flag')} This function controls whether or not smoothing
is applied to the $B$ and $H$ fields, which are naturally
piece-wise constant over each element.  Setting {\tt flag} equal to
{\tt 'on'} turns on smoothing, and setting {\tt flag} to {\tt
'off'} turns off smoothing.
\item{\tt mo\_setweightingscheme(n)} changes the weighting method used for Weighted Stress Tensor force and torque integrals. Type 0, the default, is recommended: it generates a line of elements in the middle of the air gap between components, minimizing error. Type 1 weights all elements equally and uses all elements in the gap. Types 2, 3 and 4 select alternative element weightings and mask transition profiles.

\item{\tt mo\_smoothon} turns on smoothing of the $B$ and $H$ fields, which are naturally piece-wise constant over each element. Equivalent to {\tt mo\_smooth('on')}.

\item{\tt mo\_smoothoff} turns off smoothing. Equivalent to {\tt mo\_smooth('off')}.

\item{\verb+mo_showgrid+} Show the grid points.
\item{\verb+mo_hidegrid+} Hide the grid points points.
\item{\verb+mo_gridsnap('flag')+}
Setting {\tt flag} to 'on' turns on snap to grid, setting {\tt
flag} to {\tt 'off'} turns off snap to grid.
\item{\tt mo\_snapgrid('flag')} setting {\tt flag} to {\tt 'on'} turns on snap to grid, setting it to {\tt 'off'} turns it off.

\item{\tt mo\_snapgridon} turns on snap-to-grid. Equivalent to {\tt mo\_snapgrid('on')}.

\item{\tt mo\_snapgridoff} turns off snap-to-grid. Equivalent to {\tt mo\_snapgrid('off')}.

\item{\verb+mo_setgrid(density,'type')+} Change the grid spacing.  The {\tt density}
parameter specifies the space between grid points, and the {\tt
type} parameter is set to {\tt 'cart'} for cartesian coordinates or
{\tt 'polar'} for polar coordinates.
\item{\verb+mo_hidedensityplot+} hides the flux density plot.
\item{\verb+mo_showdensityplot(legend,gscale,upper_B,lower_B,type)+}
Shows the flux density plot with options:
        \begin{itemize}
        \item {\tt legend} Set to {\tt 0} to hide the plot legend or {\tt 1} to show the plot
        legend.
        \item {\tt gscale} Set to {\tt 0} for a colour density plot or {\tt 1} for a grey scale density
        plot.
        \item{\verb+upper_B+} Sets the upper display limit for the density
        plot.
        \item{\verb+lower_B+} Sets the lower display limit for the density
        plot.
        \item{\tt type} Type of density plot to display. Valid
        entries are {\tt 'bmag'}, {\tt 'breal'}, {\tt 'bimag'}, and {\tt "logb"} for
        magnitude, real component, and imaginary component, and log magnitude of $B$;
        and {\tt "hmag"}, {\tt "hreal"}, and {\tt "himag"} for
        magnitude, real component, and imaginary component of field intensity ($H$),
        respectively.  Alternatively, current density can be
        displayed by specifying {\tt 'jmag'}, {\tt 'jreal'}, and {\tt 'jimag'} for
        magnitude, real component, and imaginary component of $J$,
        respectively.
        \end{itemize}
\item{\verb+mo_hidecontourplot+} Hides the contour plot.
\item{\verb+mo_showcontourplot(numcontours,lower_A,upper_A,type)+}
shows the $A$ contour plot with options:
        \begin{itemize}
        \item{\tt numcontours} Number of $A$ equipotential lines
        to be plotted.
        \item{\verb+upper_A+} Upper limit for $A$ contours.
        \item{\verb+lower_A+} Lower limit for $A$ contours.
                \item{\verb+type+} Choice of {\tt 'real'}, {\tt 'imag'}, or {\tt 'both'}
                        to show either the real, imaginary of both real and imaginary components of A.
        \end{itemize}

\item{\tt mo\_showvectorplot(type,scalefactor)}
controls the display of vectors denoting the field strength and
direction. The parameters taken are the \texttt{type} of plot,
which should be set to 0 for no vector plot,  1 for the real part of flux density B;
2 for the real part of field intensity H; 3 for the imaginary part of B;
4 for the imaginary part of H; 5 for both the real and imaginary parts of B;
and 6 for both the real and imaginary parts of H. The \texttt{scalefactor}
determines the relative length of the vectors. If the scale is set
to 1, the length of the vectors are chosen so that the highest flux
density corresponds to a vector that is the same length as the
current grid size setting.

\item{\tt mo\_minimize} minimizes the active magnetics output view.
\item{\tt mo\_maximize} maximizes the active magnetics output view.
\item{\tt mo\_restore} restores the active magnetics output view from a
 minimized or maximized state.
\item{\tt mo\_resize(width,height)} resizes the active magnetics output
 window client area to width $\times$ height.
\item{\tt mo\_getview} grabs the currently displayed magnetics output view and
 returns it as an image matrix, in the form \texttt{imread} would return.
\end{itemize}

\subsection{Miscellaneous}
\begin{itemize}

\item{\tt mo\_close} Closes the current post-processor instance.
\item{\tt mo\_refreshview} Redraws the current view.
\item{\tt mo\_reload} Reloads the solution from disk.
\item{\tt mo\_savebitmap('filename')} saves a bitmapped screen shot of the current
view to the file specified by {\tt 'filename'}.
\item{\tt mo\_savemetafile('filename')} saves a metafile screenshot of the current
view to the file specified by {\tt 'filename'}.
\item{\tt mo\_shownames(flag)} This function allow the user to display or hide the block label
names on screen.  To hide the block label names, {\tt flag} should be 0.  To display the
names, the parameter should be set to 1.

\item{\tt mo\_hidenames} hides the block label names. Equivalent to {\tt mo\_shownames(0)}.

\item{\tt mo\_gettitle} returns the file path of the current document as a string.

\item{\tt mo\_numnodes} Returns the number of nodes in the in focus magnetics output mesh.
\item{\tt mo\_numelements} Returns the number of elements in the in focus magnets output mesh.
\item{\tt mo\_getnode(n)} Returns the (x,y) or (r,z) position of the nth mesh node.
\item{\tt mo\_getelement(n)} returns the following proprerties for the nth element: 
	\begin{enumerate}
	    \item Index of first element node
   		\item Index of second element node
   	 	\item Index of third element node
    	\item x (or r) coordinate of the element centroid
    	\item y (or z) coordinate of the element centroid
    	\item element area using the length unit defined for the problem
    	\item group number associated with the element
	\end{enumerate}


\item {\bf Bulk getters.}  These fetch the whole mesh in one call instead of one item at a time, which is much faster than looping the per-item getters.  Mesh node indices are {\bf 1-based}, matching {\tt mo\_getnode} and {\tt mo\_getelement}.  Each returns:

\begin{itemize}
\item{\tt mo\_getallnodes} an $N \times 2$ matrix of the coordinates of every mesh node, one row per node.
\item{\tt mo\_getallelements} an $N \times 7$ matrix with one row per mesh element, holding the same seven fields as {\tt mo\_getelement}: the three node indices, the centroid coordinates, the area and the group number.
\item{\tt mo\_getallb} an $N \times 2$ matrix of the flux density at every element centroid. Complex at a nonzero frequency.
\item{\tt mo\_getalla} the vector potential at every element centroid. Complex at a nonzero frequency.
\item{\tt mo\_getallgapb('BdryName')} an $npts \times 2$ matrix of $B_r$ and $B_\theta$ at equally spaced angles around the named air gap element. An optional second argument sets $npts$ (default 360).
\item{\tt mo\_getallgapa('BdryName')} the vector potential at $npts$ points around the named air gap element. An optional second argument sets $npts$.
\item{\tt mo\_getallgapharmonics('BdryName')} an $N \times 6$ matrix of the air gap element's harmonic coefficients.
\end{itemize}

\end{itemize}


% Electrostatics-specific functions
\section{Electrostatics Preprocessor Command Set}

A number of different commands are available in the preprocessor.
Two naming conventions can be used: one which separates words in
the command names by underscores, and one that eliminates the
underscores.

\subsection{Object Add/Remove Commands}
\begin{itemize}
\item {\tt ei\_addnode(x,y)} Add a new node at x,y

\item {\tt ei\_addsegment(x1,y1,x2,y2)} Add a new line segment from node closest to
(x1,y1) to node closest to (x2,y2)

\item {\tt ei\_addblocklabel(x,y)} Add a new block label at (x,y)

\item {\tt ei\_addarc(x1,y1,x2,y2,angle,maxseg)} Add a new arc segment from the
nearest node to (x1,y1) to the nearest node to (x2,y2) with angle `angle'
divided into `maxseg' segments.

\item{\tt ei\_drawline(x1,y1,x2,y2)} Adds nodes at (x1,y1) and
(x2,y2) and adds a line between the nodes.

\item{\tt ei\_drawpolyline([x1,y1;x2,y2'...])} Adds nodes at each
of the specified points and connects them with segments.

\item{\tt ei\_drawpolygon([x1,y1;x2,y2'...])} Adds nodes at each
of the specified points and connects them with segments to form a
closed contour.

\item{\tt ei\_drawarc(x1,y1,x2,y2,angle,maxseg)} Adds nodes at
(x1,y1) and (x2,y2) and adds an arc of the specified angle and
discretization connecting the nodes.

\item{\tt ei\_drawrectangle(x1,y1,x2,y2)} Adds nodes at the
corners of a rectangle defined by the points (x1,y1) and (x2,y2),
then adds segments connecting the corners of the rectangle.

\item {\tt ei\_deleteselected} Delete all selected objects.

\item {\tt ei\_deleteselectednodes} Delete selected nodes.

\item {\tt ei\_deleteselectedlabels} Delete selected block labels.

\item {\tt ei\_deleteselectedsegments} Delete selected segments.

\item {\tt ei\_deleteselectedarcsegments} Delete selects arcs.
\end{itemize}


\subsection{Geometry Selection Commands}

\begin{itemize}
\item {\tt ei\_clearselected} Clear all selected nodes, blocks, segments and arc
segments.

\item {\tt ei\_selectsegment(x,y)} Select the line segment closest to (x,y)

\item {\tt ei\_selectnode(x,y)} Select the node closest to (x,y).
Returns the coordinates of the selected node.

\item {\tt ei\_selectlabel(x,y)} Select the label closet to (x,y).
Returns the coordinates of the selected label.

\item {\tt ei\_selectarcsegment(x,y)} Select the arc segment closest to (x,y)

\item {\tt ei\_selectgroup(n)} Select the n$^{th}$ group of nodes, segments, arc
segments and block labels. This function will clear all previously selected
elements and leave the edit mode in 4 (group)

\item{\tt ei\_selectcircle(x,y,R,editmode)} selects objects within a circle of radius
R centered at (x,y).  If only x, y, and R parameters are given, the current
edit mode is used.  If the editmode parameter is used, 0 denotes nodes, 1 denotes segments,
2 denotes block labels, 3 denotes arcs, and 4 specifies that all entity types are to be selected.

\item{\tt ei\_selectrectangle(x1,y1,x2,y2,editmode)} selects objects within a rectangle
defined by points (x1,y1) and (x2,y2). If no editmode parameter is supplied,
the current edit mode is used.  If the editmode parameter is used, 0 denotes nodes, 1 denotes segments,
2 denotes block labels, 3 denotes arcs, and 4 specifies that all entity types are to be selected.

\end{itemize}


\subsection{Object Labeling Commands}

\begin{itemize}
\item {\tt ei\_setnodeprop('propname',groupno, 'inconductor')} Set the selected
nodes to have the nodal property \texttt{'propname'} and group
number \texttt{groupno}. The \texttt{'inconductor'} string
specifies which conductor the node belongs to. If the node doesn't
belong to a named conductor, this parameter can be set to
\texttt{'<None>'}.

\item {\tt ei\_setblockprop('blockname', automesh, meshsize, group)} Set the
selected block labels to have the properties:

Block property \texttt{'blockname'}.

\texttt{automesh}: 0 = mesher defers to mesh size constraint defined in
\texttt{meshsize}, 1 = mesher automatically chooses the mesh density.

\texttt{meshsize}: size constraint on the mesh in the block marked by this
label.

A member of group number \texttt{group}

\item {\tt ei\_setsegmentprop('name', elmsize, automesh, hide, group,
'inconductor')}\\ Set the select segments to have:

Boundary property \texttt{'name'}

Local element size along segment no greater than
\texttt{elmsize}

\texttt{automesh}: 0 = mesher defers to the element constraint defined by
\texttt{elementsize}, 1 = mesher automatically chooses mesh size along the
selected segments

\texttt{hide}: 0 = not hidden in post-processor, 1 == hidden in post
processor

A member of group number \texttt{group}

A member of the conductor specified by the string
\texttt{'inconductor'.} If the segment is not part of a conductor,
this parameter can be specified as
\texttt{'<None>'}.

\item {\tt ei\_setarcsegmentprop(maxsegdeg, 'propname', hide, group,
'inconductor')} Set the selected arc segments to:

Meshed with elements that span at most \texttt{maxsegdeg} degrees
per element

Boundary property \texttt{'propname'}

\texttt{hide}: 0 = not hidden in post-processor, 1 == hidden in post
processor

A member of group number \texttt{group}

A member of the conductor specified by the string
\texttt{'inconductor'.} If the segment is not part of a conductor,
this parameter can be specified as
\texttt{'<None>'}.

\item{\tt ei\_setgroup(n)} Set the group associated of the selected items to n.

\end{itemize}

\subsection{Problem Commands}

\begin{itemize}
\item {\tt ei\_probdef(units,type,precision,depth,minangle,prevsoln,timestep)} changes the problem
definition. The \texttt{units} parameter specifies the units used
for measuring length in the problem domain. Valid \texttt{'units'}
entries are
\texttt{'inches'}, \texttt{'millimeters'}, \texttt{'centimeters'},
\texttt{'mils'}, \texttt{'meters}, and \texttt{'micrometers'}. Set
\texttt{problemtype} to \texttt{'planar'} for a 2-D planar problem, or to
\texttt{'axi'} for an axisymmetric problem. The \texttt{precision} parameter
dictates the precision required by the solver. For example, entering
\texttt{1.E-8} requires the RMS of the residual to be less than 10$^{ - 8}$.
The fourth parameter represents the depth of the problem in the
into-the-page direction for 2-D planar problems--just enter 0 as a placeholder for axisymmetric problems.
The fifth parameter represents the minimum angle constraint sent to the mesh generator (usually 30 degrees).
The sixth parameter indicates the solution from the previous time step
to be used in time-transient problems. The seventh parameter is the time step assumed
for time transient problems.

\item {\tt ei\_analyze(flag)} runs the electrostatics solver.
The {\tt flag} parameter
controls whether the solver window is visible or minimized.  For a
visible window,  specify {\tt 0}. For a minimized window, {\tt flag} should be set to {\tt 1}.
If no value is specified for {\tt flag}, the visibility of the solver is inherited from the main window,
{\em i.e.} if the main window is minimized, the solver runs minimized, too.

\item {\tt ei\_loadsolution} loads and displays the solution corresponding to the
current geometry.

\item {\tt ei\_setfocus('documentname')} Switches the
electrostatics input file upon which commands are to act. If
more than one electrostatics input file is being edited at a time,
this command can be used to switch between files so that the
mutiple files can be operated upon programmatically. {\tt
'documentname'} should contain the name of the desired document as
it appears on the window's title bar.

\item {\tt ei\_saveas('filename')} saves the file with name
\texttt{'filename'}.

\end{itemize}

\subsection{Mesh Commands}

\begin{itemize}
\item \texttt{ei\_createmesh} runs triangle to create a mesh. Note that this is not
a necessary precursor of performing an analysis, as
\texttt{ei\_analyze} will make sure the mesh is up to date before
running an analysis. Returns the number of elements in the mesh.

\item \texttt{ei\_showmesh} toggles the flag that shows or hides the mesh.

\item \texttt{ei\_purgemesh} clears the mesh out of both the screen and memory.
\end{itemize}

\subsection{Editing Commands}
\begin{itemize}
\item{\tt ei\_copyrotate(bx, by, angle, copies )}
        \begin{itemize}
        \item{\tt bx, by} -- base point for rotation
        \item{\tt angle} -- angle by which the selected objects are incrementally
        shifted to make each copy.  {\tt angle} is measured in degrees.
        \item{\tt copies} -- number of copies to be produced from
        the selected objects.
        \end{itemize}

\item{\tt ei\_copyrotate2(bx, by, angle, copies, editaction )}
        \begin{itemize}
        \item{\tt bx, by} -- base point for rotation
        \item{\tt angle} -- angle by which the selected objects are incrementally
        shifted to make each copy.  {\tt angle} is measured in degrees.
        \item{\tt copies} -- number of copies to be produced from
        the selected objects.
        \item{\tt editaction}  0 --nodes, 1 -- lines (segments), 2 --block labels, 3 -- arc
                segments, 4- group
        \end{itemize}

\item{\tt ei\_copytranslate(dx, dy, copies)}
        \begin{itemize}
        \item{\tt dx,dy} -- distance by which the selected objects are incrementally shifted.
        \item{\tt copies} -- number of copies to be produced from the selected objects.
        \end{itemize}

\item{\tt ei\_copytranslate2(dx, dy, copies, editaction)}
        \begin{itemize}
        \item{\tt dx,dy} -- distance by which the selected objects are incrementally shifted.
        \item{\tt copies} -- number of copies to be produced from the selected objects.
        \item{\tt editaction}  0 --nodes, 1 -- lines (segments), 2 --block labels, 3 -- arc
                segments, 4- group
        \end{itemize}

\item{\tt ei\_createradius(x, y, r)} turns a corner located at {\tt (x,y)} into a curve of radius {\tt r}.

\item{\tt ei\_moverotate(bx,by,shiftangle)}
        \begin{itemize}
        \item{\tt bx, by} -- base point for rotation
        \item{\tt shiftangle} -- angle in degrees by which the selected objects are rotated.
        \end{itemize}

\item{\tt ei\_moverotate2(bx,by,shiftangle, editaction)}
        \begin{itemize}
        \item{\tt bx, by} -- base point for rotation
        \item{\tt shiftangle} -- angle in degrees by which the selected objects are rotated.
        \item{\tt editaction}  0 --nodes, 1 -- lines (segments), 2 --block labels, 3 -- arc
                segments, 4- group
        \end{itemize}

\item{\tt ei\_movetranslate(dx,dy)}
        \begin{itemize}
        \item{\tt dx,dy} -- distance by which the selected objects are shifted.
        \end{itemize}

\item{\tt ei\_movetranslate2(dx,dy,editaction)}
        \begin{itemize}
        \item{\tt dx,dy} -- distance by which the selected objects are shifted.
        \item{\tt editaction}  0 --nodes, 1 -- lines (segments), 2 --block labels, 3 -- arc
                segments, 4- group
        \end{itemize}

\item{\tt ei\_scale(bx,by,scalefactor)}
        \begin{itemize}
        \item{\tt bx, by} -- base point for scaling
        \item{\tt scalefactor} -- a multiplier that determines how
        much the selected objects are scaled
        \end{itemize}

\item{\tt ei\_scale2(bx,by,scalefactor,editaction)}
        \begin{itemize}
        \item{\tt bx, by} -- base point for scaling
        \item{\tt scalefactor} -- a multiplier that determines how
        much the selected objects are scaled
        \item{\tt editaction}  0 --nodes, 1 -- lines (segments), 2 --block labels, 3 -- arc
                segments, 4- group
        \end{itemize}


\item{\tt ei\_mirror(x1,y1,x2,y2)}
mirror the selected objects about a line passing through the points
{\tt (x1,y1)} and {\tt (x2,y2)}.

\item{\tt ei\_mirror2(x1,y1,x2,y2,editaction)}
mirror the selected objects about a line passing through the points
{\tt (x1,y1)} and {\tt (x2,y2)}. Valid {\tt editaction} entries are
0 for nodes, 1 for lines (segments), 2 for block labels, 3 for arc
segments, and 4 for groups.

\item{\tt ei\_seteditmode(editmode)}
Sets the current editmode to:
        \begin{itemize}
        \item{\tt 'nodes'} - nodes
        \item{\tt 'segments'} - line segments
        \item{\tt 'arcsegments'} - arc segments
        \item{\tt 'blocks'} - block labels
        \item{\tt 'group'} - selected group
        \end{itemize}
This command will affect all subsequent uses of the other editing
commands, if they are used WITHOUT the {\tt editaction} parameter.
\end{itemize}

\subsection{Zoom Commands}

\begin{itemize}
\item \texttt{ei\_zoomnatural} zooms to a ``natural'' view with sensible extents.

\item \texttt{ei\_zoomout} zooms out by a factor of 50{\%}.

\item \texttt{ei\_zoomin} zoom in by a factor of 200{\%}.

\item \texttt{ei\_zoom(x1,y1,x2,y2)} Set the display area to be from the bottom left
corner specified by \texttt{(x1,y1}) to the top right corner
specified by \texttt{(x2,y2)}.
\end{itemize}

\subsection{View Commands}
\begin{itemize}
\item \texttt{ei\_showgrid} Show the grid points.

\item \texttt{ei\_hidegrid} Hide the grid points points.

\item \texttt{ei\_gridsnap('flag')} Setting flag to ''on'' turns on snap to grid,
setting flag to 'off'turns off snap to grid.

\item{\tt ei\_snapgridon} turns on snap-to-grid. Equivalent to {\tt ei\_gridsnap('on')}.

\item{\tt ei\_snapgridoff} turns off snap-to-grid. Equivalent to {\tt ei\_gridsnap('off')}.

\item \texttt{ei\_setgrid(density,'type')} Change the grid spacing. The density
parameter specifies the space between grid points, and the type
parameter is set to \texttt{'cart'} for Cartesian coordinates or
\texttt{'polar'} for polar coordinates.

\item \texttt{ei\_refreshview} Redraws the current view.

\item{\tt ei\_minimize} minimizes the active electrostatics input view.
\item{\tt ei\_maximize} maximizes the active electrostatics input view.
\item{\tt ei\_restore} restores the active electrostatics input view from a
 minimized or maximized state.
\item{\tt ei\_resize(width,height)} resizes the active electrostatics input
 window client area to width $\times$ height.
\item{\tt ei\_getview} grabs the currently displayed electrostatics input view and
 returns it as an image matrix, in the form \texttt{imread} would return.
\end{itemize}




\subsection{Object Properties}

\begin{itemize}
\item \texttt{ei\_getmaterial('materialname')} fetches the material specified by \texttt{materialname} 
from the materials library.

\item {\tt ei\_addmaterial('matname', ex, ey, qv)} adds a new material with
called {\tt 'matname'} with the material properties:

\texttt{ex} Relative permittivity in the x- or r-direction.

\texttt{ey} Relative permittivity in the y- or z-direction.

\texttt{qv} Volume charge density in units of C/m$^{3}$

\item \texttt{ei\_addpointprop('pointname',Vp,qp)} adds a new point property of
name \texttt{'pointname'} with either a specified potential
\texttt{Vp }a point charge density \texttt{qp} in units of C/m. .

\item \texttt{ei\_addboundprop('boundname', Vs, qs, c0, c1, BdryFormat)} adds a
new boundary property with name \texttt{'boundname'}

For a ``Fixed Voltage'' type boundary condition, set the
\texttt{Vs} parameter to the desired voltage and all other
parameters to zero.

To obtain a ``Mixed'' type boundary condition, set \texttt{C1} and
\texttt{C0} as required and \texttt{BdryFormat} to 1. Set all other
parameters to zero.

To obtain a prescribes surface charge density, set \texttt{qs} to
the desired charge density in C/m$^{2}$ and set \texttt{BdryFormat}
to 2.

For a ``Periodic'' boundary condition, set \texttt{BdryFormat} to 3
and set all other parameters to zero.

For an ``Anti-Perodic'' boundary condition, set \texttt{BdryFormat}
to 4 set all other parameters to zero.

\item \texttt{ei\_addconductorprop('conductorname', Vc, qc, conductortype)} adds a new
conductor property with name \texttt{'conductorname'} with either a
prescribed voltage or a prescribed total charge. Set the unused
property to zero. The \texttt{conductortype} parameter is 0 for
prescribed charge and 1 for prescribed voltage.

\item \texttt{ei\_deletematerial('materialname')} deletes the material named
\texttt{'materialname'}.

\item {\tt ei\_deleteboundprop('boundname')} deletes the boundary property
named {\tt 'boundname'}.

\item \texttt{ei\_deleteconductor('conductorname')} deletes the conductor named
\texttt{conductorname}.

\item \texttt{ei\_deletepointprop('pointname')} deletes the point property named
\texttt{'pointname'}

\item \texttt{ei\_modifymaterial('BlockName',propnum,value)} This function allows for
modification of a material's properties without redefining the
entire material (e.g. so that current can be modified from run to
run). The material to be modified is specified by
\texttt{'BlockName'}. The next parameter is the number of the
property to be set. The last number is the value to be applied to
the specified property. The various properties that can be modified
are listed below:

\begin{tabular}{lll}
propnum & Symbol &  Description \\ \hline
 0 & \texttt{BlockName} & Name of the material \\
 1 & \texttt{ex} & x (or r) direction relative permittivity \\
 2 & \texttt{ey} & y (or z) direction relative permittivity \\
 3 & \texttt{qs} & Volume charge
\end{tabular}


\item \texttt{ei\_modifyboundprop('BdryName',propnum,value) }This function allows for
modification of a boundary property. The BC to be modified is specified by
\texttt{'BdryName'}. The next parameter is the number of the property to be
set. The last number is the value to be applied to the specified property.
The various properties that can be modified are listed below:

\begin{tabular}{lll}
 \texttt{propnum} & Symbol & Description \\ \hline
 0 & \texttt{BdryName} & Name of boundary property \\
 1 & \texttt{Vs} & Fixed Voltage \\
 2 & \texttt{qs} & Prescribed charge density \\
 3 & \texttt{c0} & Mixed BC parameter \\
 4 & \texttt{c1} & Mixed BC parameter \\
 5 & \texttt{BdryFormat} & Type of boundary condition:\\
   &   &
   \begin{tabular}{lll}
   0 & = & Prescribed V \\
   1 & = & Mixed \\
   2 & = & Surface charge density \\
   3 & = & Periodic \\
   4 & = & Antiperiodic
   \end{tabular}
\end{tabular}





\item \texttt{ei\_modifypointprop('PointName',propnum,value)} This function allows for
modification of a point property. The point property to be modified
is specified by \texttt{'PointName'}. The next parameter is the
number of the property to be set. The last number is the value to
be applied to the specified property. The various properties that
can be modified are listed below:



\begin{tabular}{lll}
\texttt{propnum} & Symbol & Description \\ \hline
 0 & \texttt{PointName} & Name of the point property \\
 1 & \texttt{Vp} & Prescribed nodal voltage \\
 2 & \texttt{qp} & Point charge density in C/m
\end{tabular}

\item \texttt{ei\_modifyconductorprop('ConductorName',propnum,value)} This function
allows for modification of a conductor property. The conductor
property to be modified is specified by \texttt{'ConductorName'}.
The next parameter is the number of the property to be set. The
last number is the value to be applied to the specified property.
The various properties that can be modified are listed below:


\begin{tabular}{lll}
\texttt{propnum} & Symbol & Description \\ \hline
  0 & \texttt{ConductorName} & Name of the conductor property \\
  1 & \texttt{Vc} &  Conductor voltage \\
  2 & \texttt{qc} &  Total conductor charge \\
  3 & \texttt{ConductorType} & 0 = Prescribed charge, 1 = Prescribed voltage
\end{tabular}

\end{itemize}

\subsection{Miscellaneous}

\begin{itemize}
\item \texttt{ei\_savebitmap('filename')} saves a bitmapped screenshot of the current
view to the file specified by \texttt{'filename'}.

\item \texttt{ei\_savemetafile('filename')} saves a metafile screenshot of the current
view to the file specified by \texttt{'filename'}.

\item \texttt{ei\_refreshview} Redraws the current view.

\item \texttt{ei\_close} closes the preprocessor window and
destroys the current document.

\item \texttt{ei\_shownames(flag)} This function allow the user to display or hide the
block label names on screen. To hide the block label names,
\texttt{flag} should be 0. To display the names, the parameter
should be set to 1.

\item{\tt ei\_hidenames} hides the block label names. Equivalent to {\tt ei\_shownames(0)}.

\item{\tt ei\_gettitle} returns the file path of the current document as a string.

\item{\tt ei\_setcomment('comment')} sets the comment/description string associated with the current problem.

\item{\tt ei\_readdxf('filename')} This function imports a dxf file specified by {\tt 'filename'}.

\item{\tt ei\_savedxf('filename')} This function saves geometry information in a dxf file specified by {\tt 'filename'}.

\item{\tt ei\_defineouterspace(Zo,Ro,Ri)} defines
an axisymmetric external region to be used in conjuction with the
Kelvin Transformation method of modeling unbounded problems.  The
{\tt Zo} parameter is the z-location of the origin of the outer region,
the {\tt Ro} parameter is the radius of the outer region, and the {\tt
Ri} parameter is the radius of the inner region ({\em i.e.} the region of
interest). In the exterior region, the permeability varies as a function of
distance from the origin of the external region.  These parameters
are necessary to define the permeability variation in the external
region.

\item{\tt ei\_attachouterspace} marks all
selected block labels as members of the external region used for
modeling unbounded axisymmetric problems via the Kelvin
Transformation.

\item{\tt ei\_detachouterspace} undefines all selected block labels
as members of the external region used for modeling unbounded axisymmetric
problems via the Kelvin Transformation.

\item{\tt ei\_attachdefault()} marks the
selected block label as the default block label.  This block label
is applied to any region that has not been explicitly labeled.

\item{\tt ei\_detachdefault()} undefines the default
attribute for the selected block labels.

\item{\tt ei\_makeABC(n,R,x,y,bc)} creates a series of circular shells that emulate the
impedance of an unbounded domain (i.e. an Improvised Asymptotic Boundary
Condition).  The {\tt n} parameter contains the number of shells to be used
(should be between 1 and 10), {\tt R} is the radius of the solution domain, and
{\tt (x,y)} denotes the center of the solution domain.  The {\tt bc} parameter should
be specified as 0 for a Dirichlet outer edge or 1 for a Neumann outer edge.
If the function is called without all the parameters, the function makes up
reasonable values for the missing parameters. 

\end{itemize}

\section{Electrostatics Post Processor Command Set}

There are a number of scripting commands designed to operate in
the postprocessor. As with the preprocessor commands, these
commands can be used with either the underscore naming or with the
no-underscore naming convention.

\subsection{Data Extraction Commands}
\begin{itemize}

\item \texttt{eo\_lineintegral(type)} Calculates the line integral for the defined contour

\begin{tabular}{ll}
\texttt{type} & Integral \\ \hline
 0 & $E \cdot t$ \\
 1 & $D \cdot n$ \\
 2 & Contour length \\
 3 & Force from stress tensor \\
 4 & Torque from stress tensor
\end{tabular}

This integral a single result for field intensity, flux density, contour length, and torque.  For force, an array with
two elements is returned representing force in the $x$ and $y$ or $r$ and $z$ directions.

\item \texttt{eo\_blockintegral(type)} Calculates a block integral for the selected
blocks

\begin{tabular}{ll}
\texttt{type} & Integral \\ \hline
 0 & Stored Energy \\
 1 & Block Cross-section \\
 2 & Block Volume \\
 3 & Average $D$ over the block \\
 4 & Average $E$ over the block \\
 5 & Weighted Stress Tensor Force \\
 6 & Weighted Stress Tensor Torque
\end{tabular}

This integral a single result for energy, volume and torque. For force, $D$ and $E$, the function returns an array with
two elements is returned representing components in the $x$ and $y$ or $r$ and $z$ directions.

\item \texttt{eo\_getpointvalues(x,y)} Gets the values associated with the point at
$(x,y)$.  The function returns an array of results whose elements represent:

\begin{tabular}{ll}
\texttt{Symbol} &  Definition \\ \hline
\texttt{V} & Voltage \\
\texttt{Dx} & x- or r- direction component of displacement \\
\texttt{Dy} & y- or z- direction component of displacement \\
\texttt{Ex} & x- or r- direction component of electric field intensity \\
\texttt{Ey} & y- or z- direction component of electric field intensity \\
\texttt{ex} & x- or r- direction component of permittivity \\
\texttt{ey} & y- or z- direction component of permittivity \\
\texttt{nrg} & electric field energy density \\
\end{tabular}

\item{\tt eo\_getv(x,y)} Gets the voltage associated with the point at (x,y).

\item{\tt eo\_getd(x,y)} Gets the electric flux density associated with the point at (x,y).
The return value is a list with two element representing $D_x$ and
$D_y$ for planar problems and $D_r$ and $D_z$ for axisymmetric problems.

\item{\tt eo\_gete(x,y)} Gets the electric field intensity associated with the point at (x,y).
The return value is a list with two element representing $E_x$ and
$E_y$ for planar problems and $E_r$ and $E_z$ for axisymmetric problems.

\item{\tt eo\_getenergydensity(x,y)} Gets the electric field energy density associated with the point at (x,y).

\item{\tt eo\_getperm(x,y)} Gets the electric permittivity associated with the point at (x,y).
The return value is a list with two element representing $\epsilon_x$ and
$\epsilon_y$ for planar problems and $\epsilon_r$ and $\epsilon_z$ for axisymmetric problems.

\item \texttt{eo\_makeplot(PlotType,NumPoints,Filename,FileFormat) }Allows Octave to access
to the X-Y plot routines. If only PlotType or only PlotType and NumPoints
are specified, the command is interpreted as a request to plot the requested
plot type to the screen. If, in addition, the Filename parameter is
specified, the plot is instead written to disk to the specified file name as
an extended metafile. If the FileFormat parameter is also, the command is
instead interpreted as a command to write the data to disk to the specfied
file name, rather than display it to make a graphical plot. Valid entries
for PlotType are:

\begin{tabular}{ll}
\texttt{PlotType} &  Definition \\ \hline
 0 & \texttt{V (Voltage)} \\
 1 & \texttt{$\vert$D$\vert$ (Magnitude of flux density)} \\
 2 & \texttt{D . n (Normal flux density)} \\
 3 & \texttt{D . t (Tangential flux density)} \\
 4 & \texttt{$\vert$E$\vert$ (Magnitude of field intensity)} \\
 5 & \texttt{E . n (Normal field intensity)} \\
 6 & \texttt{E . t (Tangential field intensity)}
\end{tabular}

\vspace*{8pt}
Valid file formats are:

\begin{tabular}{ll}
\texttt{FileFormat} &  Definition  \\ \hline
 0 & \texttt{Multi-column text with legend } \\
 1 & \texttt{Multi-column text with no legend } \\
 2 & \texttt{Mathematica-style formatting}
\end{tabular}

For example, if one wanted to plot $V $ to the screen with 200 points evaluated
to make the graph, the command would be: \newline
\texttt{eo\_makeplot(0,200) } \newline
If this plot were to be written to disk as a metafile, the command
would be:\texttt{ } \newline
\texttt{eo\_makeplot(0,200,'c:temp.emf')} \newline
To write data instead of a plot to disk, the command would be of
the form:\texttt{ } \newline
\texttt{eo\_makeplot(0,200,'c:temp.txt',0)} \newline


\item \texttt{eo\_getprobleminfo }Returns info on problem description. Returns three
values: the Problem type (0 for planar and 1 for axisymmetric); the depth
assumed for planar problems in units of meters; and the length unit used to draw the problem, represented in meters.

\item \texttt{eo\_getconductorproperties('conductor')} Properties are returned for the
conductor property named 'conductor'. Two values are returned: The voltage
of the specified conductor, and the charge carried on the specified
conductor.
\end{itemize}

\subsection{Selection Commands}

\begin{itemize}
\item \texttt{eo\_seteditmode(mode) }Sets the mode of the postprocessor to point,
contour, or area mode. Valid entries for mode are\texttt{ 'point',
'contour', }and\texttt{ 'area'. }

\item \texttt{eo\_selectblock(x,y) }Select the block that contains point\texttt{
(x,y).}

\item \texttt{eo\_groupselectblock(n) }Selects all of the blocks that are labeled by
block labels that are members of group n. If no number is specified
({\em i.e.} \texttt{eo\_groupselectblock}), all blocks are selected.

\item \texttt{eo\_selectconductor('name')} Selects all nodes, segments, and arc
segments that are part of the conductor specified by the string
\texttt{('name')}. This command is used to select conductors for the
purposes of the ``weighted stress tensor'' force and torque
integrals, where the conductors are points or surfaces, rather than
regions ({\em i.e.} can't be selected with \texttt{eo\_selectblock}).

\item \texttt{eo\_addcontour(x,y) }Adds a contour point at\texttt{ (x,y). }If this is
the first point then it starts a contour, if there are existing
points the contour runs from the previous point to this point. The
\texttt{eo\_addcontour }command has the same functionality as a
right-button-click contour point addition when the program is
running in interactive mode.\texttt{ }

\item \texttt{eo\_bendcontour(angle,anglestep)} Replaces the straight line formed by
the last two points in the contour by an arc that spans angle
degrees. The arc is actually composed of many straight lines, each
of which is constrained to span no more than anglestep degrees. The
\texttt{angle} parameter can take on values from -180 to 180
degrees. The \texttt{anglestep} parameter must be greater than
zero. If there are less than two points defined in the contour,
this command is ignored.

\item \texttt{eo\_selectpoint(x,y)} Adds a contour point at the closest input point to
\texttt{(x,y).} If the selected point and a previous selected points lie at
the ends of an arcsegment, a contour is added that traces along the
arcsegment. The \texttt{selectpoint} command has the same
functionality as the left-button-click contour point selection when
the program is running in interactive mode.

\item \texttt{eo\_clearcontour} Clear a prevously defined contour

\item \texttt{eo\_clearblock} Clear block selection
\end{itemize}



\subsection{Zoom Commands}

\begin{itemize}
\item \texttt{eo\_zoomnatural} Zoom to the natural boundaries of the geometry.

\item \texttt{eo\_zoomin} Zoom in one level.

\item \texttt{eo\_zoomout} Zoom out one level.

\item \texttt{eo\_zoom(x1,y1,x2,y2)} Zoom to the window defined by lower left corner
\texttt{(x1,y1)} and upper right corner\texttt{ (x2,y2).}
\end{itemize}



\subsection{View Commands}

\begin{itemize}
\item \texttt{eo\_showmesh} Show the mesh.

\item \texttt{eo\_hidemesh} Hide the mesh.

\item \texttt{eo\_showpoints} Show the node points from the input geometry.

\item \texttt{eo\_hidepoints} Hide the node points from the input geometry.

\item \texttt{eo\_smooth('flag')} This function controls whether or not smoothing is
applied to the $D$ and $E$ fields which are naturally piece-wise
constant over each element. Setting flag equal to \texttt{'on'}
turns on smoothing, and setting flag to \texttt{'off'} turns off
smoothing.

\item{\tt eo\_smoothon} turns on smoothing of the $D$ and $E$ fields, which are naturally piece-wise constant over each element. Equivalent to {\tt eo\_smooth('on')}.

\item{\tt eo\_smoothoff} turns off smoothing. Equivalent to {\tt eo\_smooth('off')}.

\item \texttt{eo\_showgrid} Show the grid points.

\item \texttt{eo\_hidegrid} Hide the grid points points.

\texttt{eo\_gridsnap('flag')} Setting flag to ''on'' turns on snap to grid,
setting flag to 'off' turns off snap to grid.

\item \texttt{eo\_setgrid(density,'type')} Change the grid spacing. The density
parameter specifies the space between grid points, and the type
parameter is set to \texttt{'cart'} for Cartesian coordinates or
\texttt{'polar'} for polar coordinates.

\item \texttt{eo\_hidedensityplot} hides the flux density plot.

\item \texttt{eo\_showdensityplot(legend,gscale,type,upper{\_}D,lower{\_}D)} Shows the
flux density plot with options:

\texttt{legend} Set to 0 to hide the plot legend or 1 to show the plot
legend.

\texttt{gscale} Set to 0 for a colour density plot or 1 for a grey scale
density plot.

\texttt{upper{\_}D} Sets the upper display limit for the density plot.

\texttt{lower{\_}D} Sets the lower display limit for the density plot.

\texttt{type} Sets the type of density plot. A value of 0 plots voltage, 1
plots the magnitude of $D$, and 2 plots the magnitude of $E$

\item \texttt{eo\_hidecontourplot} Hides the contour plot.

\item \texttt{eo\_showcontourplot(numcontours,lower{\_}V,upper{\_}V)} shows the
$V$ contour plot with options:

\texttt{numcontours} Number of equipotential lines to be plotted.

\texttt{upper{\_}V} Upper limit for contours.

\texttt{lower{\_}V} Lower limit for contours.

\item \texttt{eo\_showvectorplot(type,scalefactor)}
controls the display of vectors denoting the field strength and
direction. The parameters taken are the \texttt{type} of plot,
which should be set to 0 for no vector plot, 1 for flux density
$D$, and 2 for field intensity $E$. The \texttt{scalefactor}
determines the relative length of the vectors. If the scale is set
to 1, the length of the vectors are chosen so that the highest flux
density corresponds to a vector that is the same length as the
current grid size setting.

\item{\tt eo\_minimize} minimizes the active electrostatics output view.
\item{\tt eo\_maximize} maximizes the active electrostatics output view.
\item{\tt eo\_restore} restores the active electrostatics output view from a
 minimized or maximized state.
\item{\tt eo\_resize(width,height)} resizes the active electrostatics output
 window client area to width $\times$ height.

\item{\tt eo\_getview} grabs the currently displayed electrostatics output view and
 returns it as an image matrix, in the form \texttt{imread} would return.

\end{itemize}

\subsection{Miscellaneous}

\begin{itemize}

\item \texttt{eo\_close} close the current postprocessor window.

\item{\tt eo\_refreshview} Redraws the current view.

\item{\tt eo\_reload} Reloads the solution from disk.

\item \texttt{eo\_savebitmap('filename')} saves a bitmapped screen shot of the current
view to the file specified by \texttt{'filename'}.

\item \texttt{eo\_savemetafile('filename')} saves a metafile screenshot of the current
view to the file specified by \texttt{'filename'}.

\item \texttt{eo\_shownames(flag)} This function allow the user to display or hide the
block label names on screen. To hide the block label names,
\texttt{flag} should be 0. To display the names, the parameter
should be set to 1.


\item{\tt eo\_hidenames} hides the block label names. Equivalent to {\tt eo\_shownames(0)}.

\item{\tt eo\_gettitle} returns the file path of the current document as a string.

\item{\tt eo\_snapgrid('flag')} setting {\tt flag} to {\tt 'on'} turns on snap to grid, setting it to {\tt 'off'} turns it off.

\item{\tt eo\_snapgridon} turns on snap-to-grid. Equivalent to {\tt eo\_snapgrid('on')}.

\item{\tt eo\_snapgridoff} turns off snap-to-grid. Equivalent to {\tt eo\_snapgrid('off')}.

\item{\tt eo\_numnodes} Returns the number of nodes in the in focus electrostatics output mesh.
\item{\tt eo\_numelements} Returns the number of elements in the in focus electrostatics output mesh.
\item{\tt eo\_getnode(n)} Returns the (x,y) or (r,z) position of the nth mesh node.
\item{\tt eo\_getelement(n)} returns the following properties for the nth element:
    \begin{enumerate}
        \item Index of first element node
        \item Index of second element node
        \item Index of third element node
        \item x (or r) coordinate of the element centroid
        \item y (or z) coordinate of the element centroid
        \item element area using the length unit defined for the problem
        \item group number associated with the element
    \end{enumerate}


\item {\bf Bulk getters.}  These fetch the whole mesh in one call instead of one item at a time, which is much faster than looping the per-item getters.  Mesh node indices are {\bf 1-based}, matching {\tt eo\_getnode} and {\tt eo\_getelement}.  Each returns:

\begin{itemize}
\item{\tt eo\_getallnodes} an $N \times 2$ matrix of the coordinates of every mesh node.
\item{\tt eo\_getallelements} an $N \times 7$ matrix with one row per mesh element, as {\tt eo\_getelement}.
\item{\tt eo\_getalld} an $N \times 2$ matrix of the electric displacement $D_x, D_y$ at every element. The value is constant over each element; no smoothing is applied and no point search is performed.
\item{\tt eo\_getallv} the voltage at every mesh node. Note that this is per {\em node}, whereas {\tt eo\_getalld} is per {\em element}.
\end{itemize}

\end{itemize}


% Heat flow-specific functions
\section{Heat Flow Preprocessor Lua Command Set}

A number of different commands are available in the preprocessor.
Two naming conventions can be used: one which separates words in
the command names by underscores, and one that eliminates the
underscores.

\subsection{Object Add/Remove Commands}
\begin{itemize}
\item {\tt hi\_addnode(x,y)} Add a new node at x,y

\item {\tt hi\_addsegment(x1,y1,x2,y2)} Add a new line segment from node closest to
(x1,y1) to node closest to (x2,y2)

\item {\tt hi\_addblocklabel(x,y)} Add a new block label at (x,y)

\item {\tt hi\_addarc(x1,y1,x2,y2,angle,maxseg)} Add a new arc segment from the
nearest node to (x1,y1) to the nearest node to (x2,y2) with angle `angle'
divided into `maxseg' segments.

\item{\tt hi\_drawline(x1,y1,x2,y2)} Adds nodes at (x1,y1) and
(x2,y2) and adds a line between the nodes.

\item{\tt hi\_drawpolyline([x1,y1;x2,y2'...])} Adds nodes at each
of the specified points and connects them with segments.

\item{\tt hi\_drawpolygon([x1,y1;x2,y2'...])} Adds nodes at each
of the specified points and connects them with segments to form a
closed contour.

\item{\tt hi\_drawarc(x1,y1,x2,y2,angle,maxseg)} Adds nodes at
(x1,y1) and (x2,y2) and adds an arc of the specified angle and
discretization connecting the nodes.

\item{\tt hi\_drawrectangle(x1,y1,x2,y2)} Adds nodes at the
corners of a rectangle defined by the points (x1,y1) and (x2,y2),
then adds segments connecting the corners of the rectangle.

\item {\tt hi\_deleteselected} Delete all selected objects.

\item {\tt hi\_deleteselectednodes} Delete selected nodes.

\item {\tt hi\_deleteselectedlabels} Delete selected block labels.

\item {\tt hi\_deleteselectedsegments} Delete selected segments.

\item {\tt hi\_deleteselectedarcsegments} Delete selects arcs.
\end{itemize}


\subsection{Geometry Selection Commands}

\begin{itemize}
\item {\tt hi\_clearselected()} Clear all selected nodes, blocks, segments and arc
segments.

\item {\tt hi\_selectsegment(x,y)} Select the line segment closest to (x,y)

\item {\tt hi\_selectnode(x,y)} Select the node closest to (x,y).
Returns the coordinates of the selected node.

\item {\tt hi\_selectlabel(x,y)} Select the label closet to (x,y).
Returns the coordinates of the selected label.

\item {\tt hi\_selectarcsegment(x,y)} Select the arc segment closest to (x,y)

\item {\tt hi\_selectgroup(n)} Select the n$^{th}$ group of nodes, segments, arc
segments and block labels. This function will clear all previously selected
elements and leave the edit mode in 4 (group)

\item{\tt hi\_selectcircle(x,y,R,editmode)} selects objects within a circle of radius
R centered at (x,y).  If only x, y, and R paramters are given, the current
edit mode is used.  If the editmode parameter is used, 0 denotes nodes, 1 denotes segments,
2 denotes block labels, 3 denotes arcs, and 4 specifies that all entity types are to be selected.

\item{\tt hi\_selectrectangle(x1,y1,x2,y2,editmode)} selects objects within a rectangle
defined by points (x1,y1) and (x2,y2). If no editmode parameter is supplied,
the current edit mode is used.  If the editmode parameter is used, 0 denotes nodes, 1 denotes segments,
2 denotes block labels, 3 denotes arcs, and 4 specifies that all entity types are to be selected.

\end{itemize}


\subsection{Object Labeling Commands}

\begin{itemize}
\item {\tt hi\_setnodeprop("propname",groupno, "inconductor")} Set the selected
nodes to have the nodal property \texttt{"propname"} and group
number \texttt{groupno}. The \texttt{"inconductor"} string
specifies which conductor the node belongs to. If the node doesn't
belong to a named conductor, this parameter can be set to
\texttt{"<None>"}.

\item {\tt hi\_setblockprop("blockname", automesh, meshsize, group)} Set the
selected block labels to have the properties:

Block property \texttt{"blockname"}.

\texttt{automesh}: 0 = mesher defers to mesh size constraint defined in
\texttt{meshsize}, 1 = mesher automatically chooses the mesh density.

\texttt{meshsize}: size constraint on the mesh in the block marked by this
label.

A member of group number \texttt{group}

\item {\tt hi\_setsegmentprop("propname", elementsize, automesh, hide, group,
"inconductor")} Set the select segments to have:

Boundary property \texttt{"propname"}

Local element size along segment no greater than
\texttt{elementsize}

\texttt{automesh}: 0 = mesher defers to the element constraint defined by
\texttt{elementsize}, 1 = mesher automatically chooses mesh size along the
selected segments

\texttt{hide}: 0 = not hidden in post-processor, 1 == hidden in post
processor

A member of group number \texttt{group}

A member of the conductor specified by the string
\texttt{"inconductor".} If the segment is not part of a conductor,
this parameter can be specified as
\texttt{"<None>"}.

\item {\tt hi\_setarcsegmentprop(maxsegdeg, "propname", hide, group,
"inconductor")} Set the selected arc segments to:

Meshed with elements that span at most \texttt{maxsegdeg} degrees
per element

Boundary property \texttt{"propname"}

\texttt{hide}: 0 = not hidden in post-processor, 1 == hidden in post
processor

A member of group number \texttt{group}

A member of the conductor specified by the string
\texttt{"inconductor".} If the segment is not part of a conductor,
this parameter can be specified as
\texttt{"<None>"}.

\item{\tt hi\_setgroup(n)} Set the group associated of the selected items to n.

\end{itemize}

\subsection{Problem Commands}

\begin{itemize}
\item {\tt hi\_probdef(units,type,precision,(depth),(minangle))} changes the problem
definition. The \texttt{units} parameter specifies the units used
for measuring length in the problem domain. Valid \texttt{"units"}
entries are
\texttt{"inches"}, \texttt{"millimeters"}, \texttt{"centimeters"},
\texttt{"mils"}, \texttt{"meters}, and \texttt{"micrometers"}. Set
\texttt{problemtype} to \texttt{"planar"} for a 2-D planar problem, or to
\texttt{"axi"} for an axisymmetric problem. The \texttt{precision} parameter
dictates the precision required by the solver. For example, entering
\texttt{1.E-8} requires the RMS of the residual to be less than 10$^{ - 8}$.
A fourth parameter, representing the depth of the problem in the
into-the-page direction for 2-D planar problems, can also be specified
for planar problems. A sixth parameter represents the minimum angle constraint sent to the mesh generator.

\item {\tt hi\_analyze(flag)} runs \texttt{hsolv} to solve the problem.The {\tt flag} parameter
controls whether the solver window is visible or minimized.  For a
visible window,  specify {\tt 0}. For a minimized window, {\tt flag} should be set to {\tt 1}.
If no value is specified for {\tt flag}, the visibility of the solver is inherited from the main window,
{\em i.e.} if the main window is minimized, the solver runs minimized, too.

\item {\tt hi\_loadsolution()} loads and displays the solution corresponding to the
current geometry.

\item {\tt hi\_setfocus(``documentname'')} Switches the
heat flow input file upon which Lua commands are to act. If
more than one heat flow input file is being edited at a time,
this command can be used to switch between files so that the
mutiple files can be operated upon programmatically via Lua. {\tt
documentname} should contain the name of the desired document as
it appears on the window's title bar.

\item {\tt hi\_saveas("filename")} saves the file with name
\texttt{"filename"}. Note if you use a path you must use two backslashes
{\em e.g.} \verb+c:\\temp\\myfile.feh+

\end{itemize}

\subsection{Mesh Commands}

\begin{itemize}
\item \texttt{hi\_createmesh()} runs the mesher to create a mesh. Note that this is not
a necessary precursor of performing an analysis, as
\texttt{hi\_analyze()} will make sure the mesh is up to date before
running an analysis. The number of elements in the mesh is pushed
back onto the lua stack.

\item \texttt{hi\_showmesh()} toggles the flag that shows or hides the mesh.

\item \texttt{hi\_purgemesh()} clears the mesh out of both the screen and memory.
\end{itemize}

\subsection{Editing Commands}
\begin{itemize}

\item{\tt hi\_copyrotate(bx, by, angle, copies )}
        \begin{itemize}
        \item{\tt bx, by} -- base point for rotation
        \item{\tt angle} -- angle by which the selected objects are incrementally
        shifted to make each copy.  {\tt angle} is measured in degrees.
        \item{\tt copies} -- number of copies to be produced from
        the selected objects.
        \end{itemize}

\item{\tt hi\_copyrotate2(bx, by, angle, copies, editaction )}
        \begin{itemize}
        \item{\tt bx, by} -- base point for rotation
        \item{\tt angle} -- angle by which the selected objects are incrementally
        shifted to make each copy.  {\tt angle} is measured in degrees.
        \item{\tt copies} -- number of copies to be produced from
        the selected objects.
        \item{\tt editaction}  0 --nodes, 1 -- lines (segments), 2 --block labels, 3 -- arc
                segments, 4- group
        \end{itemize}

\item{\tt hi\_copytranslate(dx, dy, copies)}
        \begin{itemize}
        \item{\tt dx,dy} -- distance by which the selected objects are incrementally shifted.
        \item{\tt copies} -- number of copies to be produced from the selected objects.
        \end{itemize}

\item{\tt hi\_copytranslate2(dx, dy, copies, editaction)}
        \begin{itemize}
        \item{\tt dx,dy} -- distance by which the selected objects are incrementally shifted.
        \item{\tt copies} -- number of copies to be produced from the selected objects.
        \item{\tt editaction}  0 --nodes, 1 -- lines (segments), 2 --block labels, 3 -- arc
                segments, 4- group
        \end{itemize}

\item{\tt hi\_createradius(x, y, r)} turns a corner located at {\tt (x,y)} into a curve of radius {\tt r}.

\item{\tt hi\_moverotate(bx,by,shiftangle)}
        \begin{itemize}
        \item{\tt bx, by} -- base point for rotation
        \item{\tt shiftangle} -- angle in degrees by which the selected objects are rotated.
        \end{itemize}

\item{\tt hi\_moverotate2(bx,by,shiftangle, editaction)}
        \begin{itemize}
        \item{\tt bx, by} -- base point for rotation
        \item{\tt shiftangle} -- angle in degrees by which the selected objects are rotated.
        \item{\tt editaction}  0 --nodes, 1 -- lines (segments), 2 --block labels, 3 -- arc
                segments, 4- group
        \end{itemize}

\item{\tt hi\_movetranslate(dx,dy)}
        \begin{itemize}
        \item{\tt dx,dy} -- distance by which the selected objects are shifted.
        \end{itemize}

\item{\tt hi\_movetranslate2(dx,dy,editaction)}
        \begin{itemize}
        \item{\tt dx,dy} -- distance by which the selected objects are shifted.
        \item{\tt editaction}  0 --nodes, 1 -- lines (segments), 2 --block labels, 3 -- arc
                segments, 4- group
        \end{itemize}

\item{\tt hi\_scale(bx,by,scalefactor)}
        \begin{itemize}
        \item{\tt bx, by} -- base point for scaling
        \item{\tt scalefactor} -- a multiplier that determines how
        much the selected objects are scaled
        \end{itemize}

\item{\tt hi\_scale2(bx,by,scalefactor,editaction)}
        \begin{itemize}
        \item{\tt bx, by} -- base point for scaling
        \item{\tt scalefactor} -- a multiplier that determines how
        much the selected objects are scaled
        \item{\tt editaction}  0 --nodes, 1 -- lines (segments), 2 --block labels, 3 -- arc
                segments, 4- group
        \end{itemize}

\item{\tt hi\_mirror(x1,y1,x2,y2)}
mirror the selected objects about a line passing through the points
{\tt (x1,y1)} and {\tt (x2,y2)}.

\item{\tt hi\_mirror2(x1,y1,x2,y2,editaction)}
mirror the selected objects about a line passing through the points
{\tt (x1,y1)} and {\tt (x2,y2)}. Valid {\tt editaction} entries are
0 for nodes, 1 for lines (segments), 2 for block labels, 3 for arc
segments, and 4 for groups.


\item \texttt{hi\_seteditmode(editmode)} Sets the current editmode to:

\texttt{"nodes"} -- nodes

\texttt{"segments"} - line segments

\texttt{"arcsegments"} - arc segments

\texttt{"blocks"} - block labels

\texttt{"group"} - selected group

This command will affect all subsequent uses of the other editing
commands, if they are used WITHOUT the \texttt{editaction}
parameter.
\end{itemize}

\subsection{Zoom Commands}

\begin{itemize}
\item \texttt{hi\_zoomnatural()} zooms to a ``natural'' view with sensible extents.

\item \texttt{hi\_zoomout()} zooms out by a factor of 50{\%}.

\item \texttt{hi\_zoomin()} zoom in by a factor of 200{\%}.

\item \texttt{hi\_zoom(x1,y1,x2,y2)} Set the display area to be from the bottom left
corner specified by \texttt{(x1,y1}) to the top right corner
specified by \texttt{(x2,y2)}.
\end{itemize}

\subsection{View Commands}
\begin{itemize}
\item \texttt{hi\_showgrid()} Show the grid points.

\item \texttt{hi\_hidegrid()} Hide the grid points points.

\item \texttt{hi\_gridsnap("flag")} Setting flag to ''on'' turns on snap to grid,
setting flag to "off"turns off snap to grid.

\item{\tt hi\_snapgridon} turns on snap-to-grid. Equivalent to {\tt hi\_gridsnap('on')}.

\item{\tt hi\_snapgridoff} turns off snap-to-grid. Equivalent to {\tt hi\_gridsnap('off')}.

\item \texttt{hi\_setgrid(density,"type")} Change the grid spacing. The density
parameter specifies the space between grid points, and the type
parameter is set to \texttt{"cart"} for Cartesian coordinates or
\texttt{"polar"} for polar coordinates.

\item \texttt{hi\_refreshview()} Redraws the current view.

\item{\tt hi\_minimize()} minimizes the active heat flow input view.

\item{\tt hi\_maximize()} maximizes the active heat flow input view.

\item{\tt hi\_restore()} restores the active heat flow input view from a
 minimized or maximized state.

\item{\tt hi\_resize(width,height)} resizes the active heat flow input
 window client area to width $\times$ height.
\item{\tt hi\_getview} grabs the currently displayed heat flow input view and
 returns it as an image matrix, in the form \texttt{imread} would return.

\end{itemize}


\subsection{Object Properties}

\begin{itemize}
\item \texttt{hi\_getmaterial('materialname')} fetches the material specified by \texttt{materialname} 
from the materials library.

\item \texttt{hi\_addmaterial("materialname", kx, ky, qv, kt)} adds a new material with
called \texttt{"materialname"} with the material properties:

\texttt{kx} Thermal conductivity in the x- or r-direction.

\texttt{ky} Thermal conductivity in the y- or z-direction.

\texttt{qv} Volume heat generation density in units of W/m$^{3}$

\texttt{kt} Volumetric heat capacity in units of MJ/(m$^{3}$*K)

\item \texttt{hi\_addpointprop("pointpropname",Tp,qp)} adds a new point property of
name \texttt{"pointpropname"} with either a specified temperature
\texttt{Tp} or a point heat generation density \texttt{qp} in units of W/m. .

\item \texttt{hi\_addboundprop("boundpropname", BdryFormat, Tset, qs, Tinf, h, beta, TinfRad)}
 adds a new boundary property with name {\tt "boundpropname"}.
\begin{itemize}
\item For a ``Fixed Temperature'' type boundary condition, set the {\tt Tset} parameter to the
desired temperature and all other parameters to zero.
\item To obtain a ``Heat Flux'' type boundary condition, set {\tt qs} to be the heat flux density and
{\tt BdryFormat} to 1. Set all other parameters to zero.
\item To obtain a convection boundary condition, set {\tt h} to the desired heat transfer coefficient and
{\tt Tinf} to the desired external temperature and set {\tt BdryFormat} to 2.  Heat flux {\tt qs} is also honored with this boundary
condition type, i.e. can model convection + flux.
\item For a Radiation boundary condition, set {\tt beta} equal to the desired emissivity and {\tt TinfRad} to the 
desired external temperature and set {\tt BdryFormat} to 3. Also includes a convection contribution and heat , so {\tt h} and {\tt Tinf} that define convection and specified heat flux {\tt qs} are also honored with this BC type.  If no {\tt TinfRad} is specified, it is assumed that {\tt Tinf} applies to both convection and radiation portions.
\item For a ``Periodic'' boundary condition, set {\tt BdryFormat} to 4 and set all other
parameters to zero.
\item For an ``Anti-Perodic'' boundary condition, set {\tt BdryFormat} to 5 set all other
parameters to zero.
\end{itemize}

\item \texttt{hi\_addconductorprop("conductorname", Tc, qc, conductortype)} adds a new
conductor property with name \texttt{"conductorname"} with either a
prescribed temperature or a prescribed total heat flux. Set the unused
property to zero. The \texttt{conductortype} parameter is 0 for
prescribed heat flux and 1 for prescribed temperature.

\item \texttt{hi\_deletematerial("materialname")} deletes the material named
\texttt{"materialname"}.

\item \texttt{hi\_deleteboundprop("boundpropname")} deletes the boundary property
named \texttt{"boundpropname"}.

\item \texttt{hi\_deleteconductor("conductorname")} deletes the conductor named
\texttt{conductorname}.

\item \texttt{hi\_deletepointprop("pointpropname")} deletes the point property named
\texttt{"pointpropname"}

\item \texttt{hi\_modifymaterial("BlockName",propnum,value)} This function allows for
modification of a material's properties without redefining the
entire material (e.g. so that current can be modified from run to
run). The material to be modified is specified by
\texttt{"BlockName"}. The next parameter is the number of the
property to be set. The last number is the value to be applied to
the specified property. The various properties that can be modified
are listed below:

\begin{tabular}{lll}
propnum & Symbol &  Description \\ \hline
 0 & \texttt{BlockName} & Name of the material \\
 1 & \texttt{kx} & x (or r) direction thermal conductivity \\
 2 & \texttt{ky} & y (or z) direction thermal conductivity \\
 3 & \texttt{qs} & Volume heat generation \\
 4 & \texttt{kt} & Volumetric heat capacity 
\end{tabular}


\item \texttt{hi\_modifyboundprop("BdryName",propnum,value) }This function allows for
modification of a boundary property. The BC to be modified is specified by
\texttt{"BdryName"}. The next parameter is the number of the property to be
set. The last number is the value to be applied to the specified property.
The various properties that can be modified are listed below:

\begin{tabular}{lll}
 \texttt{propnum} & Symbol & Description \\ \hline
 0 & \texttt{BdryName} & Name of boundary property \\ 
 1 & \texttt{BdryFormat} & Type of boundary condition:\\
   &   &
   \begin{tabular}{lll}
   0 & = & Prescribed temperature \\
   1 & = & Heat Flux \\
   2 & = & Convection \\
   3 & = & Radiation \\
   4 & = & Periodic \\
   5 & = & Antiperiodic
   \end{tabular} \\
 2 & \texttt{Tset} & Fixed Temperature \\
 3 & \texttt{qs} & Prescribed heat flux density \\
 4 & \texttt{Tinf} & External temperature \\
 5 & \texttt{h} & Heat transfer coefficient \\
 6 & \texttt{beta} & Emissivity \\
 7 & \texttt{TinfRad} & External temperature for radiation \\
\end{tabular}

\item \texttt{hi\_modifypointprop("PointName",propnum,value)} This function allows for
modification of a point property. The point property to be modified
is specified by \texttt{"PointName"}. The next parameter is the
number of the property to be set. The last number is the value to
be applied to the specified property. The various properties that
can be modified are listed below:

\begin{tabular}{lll}
\texttt{propnum} & Symbol & Description \\ \hline
 0 & \texttt{PointName} & Name of the point property \\
 1 & \texttt{Tp} & Prescribed nodal temperature \\
 2 & \texttt{qp} & Point heat generation in W/m
\end{tabular}

\item \texttt{hi\_modifyconductorprop("ConductorName",propnum,value)} This function
allows for modification of a conductor property. The conductor
property to be modified is specified by \texttt{"ConductorName"}.
The next parameter is the number of the property to be set. The
last number is the value to be applied to the specified property.
The various properties that can be modified are listed below:

\begin{tabular}{lll}
\texttt{propnum} & Symbol & Description \\ \hline
  0 & \texttt{ConductorName} & Name of the conductor property \\
  1 & \texttt{Tc} &  Conductor temperature \\
  2 & \texttt{qc} &  Total conductor heat flux \\
  3 & \texttt{ConductorType} & 0 = Prescribed heat flow, 1 = Prescribed temperature
\end{tabular}

\item \texttt{hi\_addtkpoint("materialname",T,k)} adds
the point {\tt (T,k)} to the thermal conductivity vs. temperature curve for the material specified by
{\tt "materialname"}.

\item{\tt hi\_addtkpoints('blockname',points)} adds a whole thermal-conductivity curve to the material named {\tt 'blockname'}, where {\tt points} is an $n$-by-2 matrix of temperature and conductivity pairs. The three-argument form takes the two columns separately. Calls {\tt hi\_addtkpoint} once per row.

\item \texttt{hi\_cleartkpoints("materialname")} erases
all of the thermal conductivity points that have been defined for the material named
{\tt "materialname"}.

\end{itemize}

\subsection{Miscellaneous}

\begin{itemize}
\item \texttt{hi\_savebitmap("filename")} saves a bitmapped screenshot of the current
view to the file specified by \texttt{"filename"}, subject to the
\texttt{printf}-type formatting explained previously for the
\texttt{hi\_saveas} command.

\item \texttt{hi\_savemetafile("filename")} saves a metafile screenshot of the current
view to the file specified by \texttt{"filename"}, subject to the
\texttt{printf}-type formatting explained previously for the
\texttt{hi\_saveas} command.

\item \texttt{hi\_refreshview()} Redraws the current view.

\item \texttt{hi\_close()} closes the preprocessor window and
destroys the current document.

\item \texttt{hi\_shownames(flag)} This function allow the user to display or hide the
block label names on screen. To hide the block label names,
\texttt{flag} should be 0. To display the names, the parameter
should be set to 1.

\item{\tt hi\_hidenames} hides the block label names. Equivalent to {\tt hi\_shownames(0)}.

\item{\tt hi\_gettitle} returns the file path of the current document as a string.

\item{\tt hi\_setcomment('comment')} sets the comment/description string associated with the current problem.

\item{\tt hi\_readdxf("filename")} This function imports a dxf file specified by {\tt "filename"}.

\item{\tt hi\_savedxf("filename")} This function saves geometry information in a dxf file specified by {\tt "filename"}.

\item{\tt hi\_defineouterspace(Zo,Ro,Ri)} defines
an axisymmetric external region to be used in conjuction with the
Kelvin Transformation method of modeling unbounded problems.  The
{\tt Zo} parameter is the z-location of the origin of the outer region,
the {\tt Ro} parameter is the radius of the outer region, and the {\tt
Ri} parameter is the radius of the inner region ({\em i.e.} the region of
interest). In the exterior region, the permeability varies as a function of
distance from the origin of the external region.  These parameters
are necessary to define the permeability variation in the external
region.

\item{\tt hi\_attachouterspace()} marks all
selected block labels as members of the external region used for
modeling unbounded axisymmetric problems via the Kelvin
Transformation.

\item{\tt hi\_detachouterspace()} undefines all selected block labels
as members of the external region used for modeling unbounded axisymmetric
problems via the Kelvin Transformation.

\item{\tt hi\_attachdefault()} marks the
selected block label as the default block label.  This block label
is applied to any region that has not been explicitly labeled.

\item{\tt hi\_detachdefault()} undefines the default
attribute for the selected block labels.

\item{\tt hi\_makeABC(n,R,x,y,bc)} creates a series of circular shells that emulate the
impedance of an unbounded domain (i.e. an Improvised Asymptotic Boundary
Condition).  The {\tt n} parameter contains the number of shells to be used
(should be between 1 and 10), {\tt R} is the radius of the solution domain, and
{\tt (x,y)} denotes the center of the solution domain.  The {\tt bc} parameter should
be specified as 0 for a Dirichlet outer edge or 1 for a Neumann outer edge.
If the function is called without all the parameters, the function makes up
reasonable values for the missing parameters. 

\end{itemize}

\section{Heat Flow Post Processor Command Set}

There are a number of Lua scripting commands designed to operate in
the postprocessor. As with the preprocessor commands, these
commands can be used with either the underscore naming or with the
no-underscore naming convention.

\subsection{Data Extraction Commands}
\begin{itemize}

\item \texttt{ho\_lineintegral(type)} Calculate the line integral for the defined contour

\begin{tabular}{ll}
\texttt{type} & Integral \\ \hline
 0 & Temperature difference ($G \cdot t$) \\
 1 & Heat flux through the contour ($F \cdot n$) \\
 2 & Contour length \\
 3 & Average Temperature \\
\end{tabular}

This integral returns either 1 or 2 values, depending on the
integral type, $e.g.:$

\texttt{Ftot, Favg = ho\_lineintegral(2)}

\item \texttt{ho\_blockintegral(type)} Calculate a block integral for the selected
blocks

\begin{tabular}{ll}
\texttt{type} & Integral \\ \hline
 0 & Average $T$ over the block \\
 1 & Block Cross-section \\
 2 & Block Volume \\
 3 & Average $F$ over the block \\
 4 & Average $G$ over the block \\
\end{tabular}

Returns one or two floating point values as results, $e.g$.:

\texttt{Gx, Gy = ho\_blockintegral(4)}

\item \texttt{ho\_getpointvalues(X,Y)} Get the values associated with the point at x,y
The return values, in order, are:

\begin{tabular}{ll}
\texttt{Symbol} &  Definition \\ \hline
\texttt{V} & Temperature \\
\texttt{Fx} & x- or r- direction component of heat flux density \\
\texttt{Fy} & y- or z- direction component of heat flux density \\
\texttt{Gx} & x- or r- direction component of temperature gradient \\
\texttt{Gy} & y- or z- direction component of temperature gradient \\
\texttt{kx} & x- or r- direction component of thermal conductivity \\
\texttt{ky} & y- or z- direction component of thermal conductivity \\
\end{tabular}

Example: To catch all values at (0.01,0) use

\texttt{T,Fx,Fy,Gx,Gy,kx,ky= ho\_getpointvalues(0.01,0) }

\item{\tt ho\_gett(x,y)} returns the temperature at {\tt (x,y)} -- column 1 of {\tt ho\_getpointvalues}. An equivalent form is {\tt ho\_gett([x,y])}.

\item{\tt ho\_getf(x,y)} returns the two components of heat flux density {\tt Fx, Fy} at {\tt (x,y)} -- columns 2 and 3 of {\tt ho\_getpointvalues}. An equivalent form is {\tt ho\_getf([x,y])}.

\item{\tt ho\_getg(x,y)} returns the two components of temperature gradient {\tt Gx, Gy} at {\tt (x,y)} -- columns 4 and 5 of {\tt ho\_getpointvalues}. An equivalent form is {\tt ho\_getg([x,y])}.

\item{\tt ho\_getk(x,y)} returns the two components of thermal conductivity {\tt kx, ky} at {\tt (x,y)} -- columns 6 and 7 of {\tt ho\_getpointvalues}. An equivalent form is {\tt ho\_getk([x,y])}.

\item \texttt{ho\_makeplot(PlotType,NumPoints,Filename,FileFormat) }Allows Lua access
to the X-Y plot routines. If only PlotType or only PlotType and NumPoints
are specified, the command is interpreted as a request to plot the requested
plot type to the screen. If, in addition, the Filename parameter is
specified, the plot is instead written to disk to the specified file name as
an extended metafile. If the FileFormat parameter is also, the command is
instead interpreted as a command to write the data to disk to the specfied
file name, rather than display it to make a graphical plot. Valid entries
for PlotType are:

\begin{tabular}{ll}
\texttt{PlotType} &  Definition \\ \hline
 0 & \texttt{V (Temperature)} \\
 1 & \texttt{$\vert$D$\vert$ (Magnitude of heat flux density)} \\
 2 & \texttt{D . n (Normal heat flux density)} \\
 3 & \texttt{D . t (Tangential heat flux density)} \\
 4 & \texttt{$\vert$E$\vert$ (Magnitude of field intensity)} \\
 5 & \texttt{E . n (Normal field intensity)} \\
 6 & \texttt{E . t (Tangential field intensity)}
\end{tabular}

\vspace*{8pt}
Valid file formats are:

\begin{tabular}{ll}
\texttt{FileFormat} &  Definition  \\ \hline
 0 & \texttt{Multi-column text with legend } \\
 1 & \texttt{Multi-column text with no legend } \\
 2 & \texttt{Mathematica-style formatting}
\end{tabular}

For example, if one wanted to plot $V $ to the screen with 200 points evaluated
to make the graph, the command would be: \newline
\texttt{ho\_makeplot(0,200) } \newline
If this plot were to be written to disk as a metafile, the command
would be:\texttt{ } \newline
\texttt{ho\_makeplot(0,200,"c:temp.emf")} \newline
To write data instead of a plot to disk, the command would be of
the form:\texttt{ } \newline
\texttt{ho\_makeplot(0,200,"c:temp.txt",0)} \newline


\item \texttt{ho\_getprobleminfo() }Returns info on problem description. Returns three
values: the Problem type (0 for planar and 1 for axisymmetric); the depth
assumed for planar problems in units of meters; and the length unit used to draw the problem, represented in meters.

\item \texttt{ho\_getconductorproperties("conductor")} Properties are returned for the
conductor property named "conductor". Two values are returned: The temperature
of the specified conductor, and the total heat flux through the specified
conductor.
\end{itemize}

\subsection{Selection Commands}

\begin{itemize}
\item \texttt{ho\_seteditmode(mode) }Sets the mode of the postprocessor to point,
contour, or area mode. Valid entries for mode are\texttt{ "point",
"contour", }and\texttt{ "area". }

\item \texttt{ho\_selectblock(x,y) }Select the block that contains point\texttt{
(x,y).}

\item \texttt{ho\_groupselectblock(n) }Selects all of the blocks that are labeled by
block labels that are members of group n. If no number is specified
({\em i.e.} \texttt{ho\_groupselectblock()}), all blocks are selected.

\item \texttt{ho\_selectconductor("name")} Selects all nodes, segments, and arc
segments that are part of the conductor specified by the string
\texttt{("name")}. This command is used to select conductors for the
purposes of the ``weighted stress tensor'' force and torque
integrals, where the conductors are points or surfaces, rather than
regions ({\em i.e.} can't be selected with \texttt{ho\_selectblock}).

\item \texttt{ho\_addcontour(x,y) }Adds a contour point at\texttt{ (x,y). }If this is
the first point then it starts a contour, if there are existing
points the contour runs from the previous point to this point. The
\texttt{ho\_addcontour }command has the same functionality as a
right-button-click contour point addition when the program is
running in interactive mode.\texttt{ }

\item \texttt{ho\_bendcontour(angle,anglestep)} Replaces the straight line formed by
the last two points in the contour by an arc that spans angle
degrees. The arc is actually composed of many straight lines, each
of which is constrained to span no more than anglestep degrees. The
\texttt{angle} parameter can take on values from -180 to 180
degrees. The \texttt{anglestep} parameter must be greater than
zero. If there are less than two points defined in the contour,
this command is ignored.

\item \texttt{ho\_selectpoint(x,y)} Adds a contour point at the closest input point to
\texttt{(x,y).} If the selected point and a previous selected points lie at
the ends of an arcsegment, a contour is added that traces along the
arcsegment. The \texttt{selectpoint} command has the same
functionality as the left-button-click contour point selection when
the program is running in interactive mode.

\item \texttt{ho\_clearcontour()} Clear a prevously defined contour

\item \texttt{ho\_clearblock()} Clear block selection
\end{itemize}



\subsection{Zoom Commands}

\begin{itemize}
\item \texttt{ho\_zoomnatural()} Zoom to the natural boundaries of the geometry.

\item \texttt{ho\_zoomin()} Zoom in one level.

\item \texttt{ho\_zoomout()} Zoom out one level.

\item \texttt{ho\_zoom(x1,y1,x2,y2)} Zoom to the window defined by lower left corner
\texttt{(x1,y1)} and upper right corner\texttt{ (x2,y2).}
\end{itemize}



\subsection{View Commands}

\begin{itemize}
\item \texttt{ho\_showmesh()} Show the mesh.

\item \texttt{ho\_hidemesh()} Hide the mesh.

\item \texttt{ho\_showpoints()} Show the node points from the input geometry.

\item \texttt{ho\_hidepoints()} Hide the node points from the input geometry.

\item \texttt{ho\_smooth("flag")} This function controls whether or not smoothing is
applied to the $F$ and $G$ fields which are naturally piece-wise
constant over each element. Setting flag equal to \texttt{"on"}
turns on smoothing, and setting flag to \texttt{"off"} turns off
smoothing.

\item{\tt ho\_smoothon} turns on smoothing of the temperature gradient and heat flux density fields, which are naturally piece-wise constant over each element. Equivalent to {\tt ho\_smooth('on')}.

\item{\tt ho\_smoothoff} turns off smoothing. Equivalent to {\tt ho\_smooth('off')}.

\item \texttt{ho\_showgrid()} Show the grid points.

\item \texttt{ho\_hidegrid()} Hide the grid points points.

\texttt{ho\_gridsnap("flag")} Setting flag to ''on'' turns on snap to grid,
setting flag to "off" turns off snap to grid.

\item \texttt{ho\_setgrid(density,"type")} Change the grid spacing. The density
parameter specifies the space between grid points, and the type
parameter is set to \texttt{"cart"} for Cartesian coordinates or
\texttt{"polar"} for polar coordinates.

\item \texttt{ho\_hidedensityplot()} hides the heat flux density plot.

\item \texttt{ho\_showdensityplot(legend,gscale,type,upper,lower)} Shows the
heat flux density plot with options:

\texttt{legend} Set to 0 to hide the plot legend or 1 to show the plot
legend.

\texttt{gscale} Set to 0 for a colour density plot or 1 for a grey scale
density plot.

\texttt{upper} Sets the upper display limit for the density plot.

\texttt{lower} Sets the lower display limit for the density plot.

\texttt{type} Sets the type of density plot. A value of 0 plots temperature, 1
plots the magnitude of $F$, and 2 plots the magnitude of $G$

\item \texttt{ho\_hidecontourplot()} Hides the contour plot.

\item \texttt{ho\_showcontourplot(numcontours,lower{\_}V,upper{\_}V)} shows the
$V$ contour plot with options:

\texttt{numcontours} Number of equipotential lines to be plotted.

\texttt{upper{\_}V} Upper limit for contours.

\texttt{lower{\_}V} Lower limit for contours.

If \texttt{ho\_numcontours} is -1 all parameters are ignored and
default values are used, \\ e.g. \texttt{show{\_}contour{\_}plot(-1)}

\item \texttt{ho\_showvectorplot(type,scalefactor)}
controls the display of vectors denoting the field strength and
direction. The parameters taken are the \texttt{type} of plot,
which should be set to 0 for no vector plot, 1 for heat flux density
$F$, and 2 for temperature gradient $G$. The \texttt{scalefactor}
determines the relative length of the vectors. If the scale is set
to 1, the length of the vectors are chosen so that the highest flux
density corresponds to a vector that is the same length as the
current grid size setting.

\item{\tt ho\_minimize()} minimizes the active heat flow input view.

\item{\tt ho\_maximize()} maximizes the active heat flow input view.

\item{\tt ho\_restore()} restores the active heat flow input view from a
 minimized or maximized state.

\item{\tt ho\_resize(width,height)} resizes the active heat flow input
 window client area to width $\times$ height.
\item{\tt ho\_getview} grabs the currently displayed heat flow output view and
 returns it as an image matrix, in the form \texttt{imread} would return.

\end{itemize}

\subsection{Miscellaneous}

\begin{itemize}

\item \texttt{ho\_close()} close the current postprocessor window.

\item{\tt ho\_refreshview()} Redraws the current view.

\item{\tt ho\_reload()} Reloads the solution from disk.

\item \texttt{ho\_savebitmap("filename")} saves a bitmapped screen shot of the current
view to the file specified by \texttt{"filename"}. Note that if you
use a path you must use two backslashes (e.g.
\texttt{"c:$\backslash \backslash $temp$\backslash \backslash
$myfile.png"}). If the file name contains a space (e.g. file
names like c:$\backslash $program files$\backslash $stuff) you must
enclose the file name in (extra) quotes by using a
\texttt{$\backslash $"} sequence. For example:

\texttt{ho\_savebitmap("$\backslash $"c:$\backslash \backslash
$temp$\backslash \backslash $screenshot.png$\backslash $"")}

\item \texttt{ho\_savemetafile("filename")} saves a metafile screenshot of the current
view to the file specified by \texttt{"filename"}, subject to the
printf-type formatting explained previously for the
\texttt{savebitmap} command.

\item \texttt{ho\_shownames(flag)} This function allow the user to display or hide the
block label names on screen. To hide the block label names,
\texttt{flag} should be 0. To display the names, the parameter
should be set to 1.


\item{\tt ho\_hidenames} hides the block label names. Equivalent to {\tt ho\_shownames(0)}.

\item{\tt ho\_gettitle} returns the file path of the current document as a string.

\item{\tt ho\_snapgrid('flag')} setting {\tt flag} to {\tt 'on'} turns on snap to grid, setting it to {\tt 'off'} turns it off.

\item{\tt ho\_snapgridon} turns on snap-to-grid. Equivalent to {\tt ho\_snapgrid('on')}.

\item{\tt ho\_snapgridoff} turns off snap-to-grid. Equivalent to {\tt ho\_snapgrid('off')}.

\item{\tt ho\_numnodes} Returns the number of nodes in the in focus heat flow output mesh.
\item{\tt ho\_numelements} Returns the number of elements in the in focus heat flow output mesh.
\item{\tt ho\_getnode(n)} Returns the (x,y) or (r,z) position of the nth mesh node.
\item{\tt ho\_getelement(n)} returns the following proprerties for the nth element:
    \begin{enumerate}
        \item Index of first element node
        \item Index of second element node
        \item Index of third element node
        \item x (or r) coordinate of the element centroid
        \item y (or z) coordinate of the element centroid
        \item element area using the length unit defined for the problem
        \item group number associated with the element
    \end{enumerate}


\item {\bf Bulk getters.}  These fetch the whole mesh in one call instead of one item at a time, which is much faster than looping the per-item getters.  Mesh node indices are {\bf 1-based}, matching {\tt ho\_getnode} and {\tt ho\_getelement}.  Each returns:

\begin{itemize}
\item{\tt ho\_getallnodes} an $N \times 2$ matrix of the coordinates of every mesh node.
\item{\tt ho\_getallelements} an $N \times 7$ matrix with one row per mesh element, as {\tt ho\_getelement}.
\item{\tt ho\_getallf} an $N \times 2$ matrix of the heat flux density $F_x, F_y$ at every element. The value is constant over each element; no smoothing is applied and no point search is performed.
\item{\tt ho\_getallt} the temperature at every mesh node. Note that this is per {\em node}, whereas {\tt ho\_getallf} is per {\em element}.
\end{itemize}

\end{itemize}


% Current flow-specific functions
\section{Current Flow Preprocessor Lua Command Set}

A number of different commands are available in the preprocessor.
Two naming conventions can be used: one which separates words in
the command names by underscores, and one that eliminates the
underscores.

\subsection{Object Add/Remove Commands}
\begin{itemize}
\item {\tt ci\_addnode(x,y)} Add a new node at x,y

\item {\tt ci\_addsegment(x1,y1,x2,y2)} Add a new line segment from node closest to
(x1,y1) to node closest to (x2,y2)

\item {\tt ci\_addblocklabel(x,y)} Add a new block label at (x,y)

\item {\tt ci\_addarc(x1,y1,x2,y2,angle,maxseg)} Add a new arc segment from the
nearest node to (x1,y1) to the nearest node to (x2,y2) with angle `angle'
divided into `maxseg' segments.

\item{\tt ci\_drawline(x1,y1,x2,y2)} Adds nodes at (x1,y1) and
(x2,y2) and adds a line between the nodes.

\item{\tt ci\_drawpolyline([x1,y1;x2,y2'...])} Adds nodes at each
of the specified points and connects them with segments.

\item{\tt ci\_drawpolygon([x1,y1;x2,y2'...])} Adds nodes at each
of the specified points and connects them with segments to form a
closed contour.

\item{\tt ci\_drawarc(x1,y1,x2,y2,angle,maxseg)} Adds nodes at
(x1,y1) and (x2,y2) and adds an arc of the specified angle and
discretization connecting the nodes.

\item{\tt ci\_drawrectangle(x1,y1,x2,y2)} Adds nodes at the
corners of a rectangle defined by the points (x1,y1) and (x2,y2),
then adds segments connecting the corners of the rectangle.

\item {\tt ci\_deleteselected} Delete all selected objects.

\item {\tt ci\_deleteselectednodes} Delete selected nodes.

\item {\tt ci\_deleteselectedlabels} Delete selected block labels.

\item {\tt ci\_deleteselectedsegments} Delete selected segments.

\item {\tt ci\_deleteselectedarcsegments} Delete selects arcs.
\end{itemize}


\subsection{Geometry Selection Commands}

\begin{itemize}
\item {\tt ci\_clearselected()} Clear all selected nodes, blocks, segments and arc
segments.

\item {\tt ci\_selectsegment(x,y)} Select the line segment closest to (x,y)

\item {\tt ci\_selectnode(x,y)} Select the node closest to (x,y).
Returns the coordinates of the selected node.

\item {\tt ci\_selectlabel(x,y)} Select the label closet to (x,y).
Returns the coordinates of the selected label.

\item {\tt ci\_selectarcsegment(x,y)} Select the arc segment closest to (x,y)

\item {\tt ci\_selectgroup(n)} Select the n$^{th}$ group of nodes, segments, arc
segments and block labels. This function will clear all previously selected
elements and leave the edit mode in 4 (group)

\item{\tt ci\_selectcircle(x,y,R,editmode)} selects objects within a circle of radius
R centered at (x,y).  If only x, y, and R paramters are given, the current
edit mode is used.  If the editmode parameter is used, 0 denotes nodes, 1 denotes segments,
2 denotes block labels, 3 denotes arcs, and 4 specifies that all entity types are to be selected.

\item{\tt ci\_selectrectangle(x1,y1,x2,y2,editmode)} selects objects within a rectangle
defined by points (x1,y1) and (x2,y2). If no editmode parameter is supplied,
the current edit mode is used.  If the editmode parameter is used, 0 denotes nodes, 1 denotes segments,
2 denotes block labels, 3 denotes arcs, and 4 specifies that all entity types are to be selected.

\end{itemize}


\subsection{Object Labeling Commands}

\begin{itemize}
\item {\tt ci\_setnodeprop("propname",groupno, "inconductor")} Set the selected
nodes to have the nodal property \texttt{"propname"} and group
number \texttt{groupno}. The \texttt{"inconductor"} string
specifies which conductor the node belongs to. If the node doesn't
belong to a named conductor, this parameter can be set to
\texttt{"<None>"}.

\item {\tt ci\_setblockprop("blockname", automesh, meshsize, group)} Set the
selected block labels to have the properties:

Block property \texttt{"blockname"}.

\texttt{automesh}: 0 = mesher defers to mesh size constraint defined in
\texttt{meshsize}, 1 = mesher automatically chooses the mesh density.

\texttt{meshsize}: size constraint on the mesh in the block marked by this
label.

A member of group number \texttt{group}

\item {\tt ci\_setsegmentprop("propname", elementsize, automesh, hide, group,
"inconductor",)} Set the select segments to have:

Boundary property \texttt{"propname"}

Local element size along segment no greater than
\texttt{elementsize}

\texttt{automesh}: 0 = mesher defers to the element constraint defined by
\texttt{elementsize}, 1 = mesher automatically chooses mesh size along the
selected segments

\texttt{hide}: 0 = not hidden in post-processor, 1 == hidden in post
processor

A member of group number \texttt{group}

A member of the conductor specified by the string
\texttt{"inconductor".} If the segment is not part of a conductor,
this parameter can be specified as
\texttt{"<None>"}.

\item {\tt ci\_setarcsegmentprop(maxsegdeg, "propname", hide, group,
"inconductor")} Set the selected arc segments to:

Meshed with elements that span at most \texttt{maxsegdeg} degrees
per element

Boundary property \texttt{"propname"}

\texttt{hide}: 0 = not hidden in post-processor, 1 == hidden in post
processor

A member of group number \texttt{group}

A member of the conductor specified by the string
\texttt{"inconductor".} If the segment is not part of a conductor,
this parameter can be specified as
\texttt{"<None>"}.

\item{\tt ci\_setgroup(n)} Set the group associated of the selected items to n.
\end{itemize}

\subsection{Problem Commands}

\begin{itemize}
\item {\tt ci\_probdef(units,type,frequency,precision,(depth),(minangle))} changes the problem
definition. The \texttt{units} parameter specifies the units used
for measuring length in the problem domain. Valid \texttt{"units"}
entries are
\texttt{"inches"}, \texttt{"millimeters"}, \texttt{"centimeters"},
\texttt{"mils"}, \texttt{"meters}, and \texttt{"micrometers"}. Set
\texttt{problemtype} to \texttt{"planar"} for a 2-D planar problem, or to
\texttt{"axi"} for an axisymmetric problem. The \texttt{frequency} parameter specifies the frequency
in Hz at which the analysis ois to be performed. The \texttt{precision} parameter
dictates the precision required by the solver. For example, entering
\texttt{1.E-8} requires the RMS of the residual to be less than 10$^{ - 8}$.
A fourth parameter, representing the depth of the problem in the
into-the-page direction for 2-D planar problems, can also be specified
for planar problems. A sixth parameter represents the minimum angle constraint sent to the mesh generator.

\item {\tt ci\_analyze(flag)} runs \texttt{belasolv} to solve the problem. 
The {\tt flag} parameter
controls whether the solver window is visible or minimized.  For a
visible window,  specify {\tt 0}. For a minimized window, {\tt flag} should be set to {\tt 1}.
If no value is specified for {\tt flag}, the visibility of the solver is inherited from the main window,
{\em i.e.} if the main window is minimized, the solver runs minimized, too.

\item {\tt ci\_loadsolution()} loads and displays the solution corresponding to the
current geometry.

\item {\tt ci\_setfocus("documentname")} Switches the
electrostatics input file upon which Lua commands are to act. If
more than one electrostatics input file is being edited at a time,
this command can be used to switch between files so that the
mutiple files can be operated upon programmatically via Lua. {\tt
documentname} should contain the name of the desired document as
it appears on the window's title bar.

\item {\tt ci\_saveas("filename")} saves the file with name
\texttt{"filename"}. Note if you use a path you must use two backslashes
{\em e.g.} \verb+c:\\temp\\myfemmfile.fee+

\end{itemize}

\subsection{Mesh Commands}

\begin{itemize}
\item \texttt{ci\_createmesh()} runs triangle to create a mesh. Note that this is not
a necessary precursor of performing an analysis, as
\texttt{ci\_analyze()} will make sure the mesh is up to date before
running an analysis. The number of elements in the mesh is pushed
back onto the lua stack.

\item \texttt{ci\_showmesh()} toggles the flag that shows or hides the mesh.

\item \texttt{ci\_purgemesh()} clears the mesh out of both the screen and memory.
\end{itemize}

\subsection{Editing Commands}
\begin{itemize}

\item{\tt ci\_copyrotate(bx, by, angle, copies )}
        \begin{itemize}
        \item{\tt bx, by} -- base point for rotation
        \item{\tt angle} -- angle by which the selected objects are incrementally
        shifted to make each copy.  {\tt angle} is measured in degrees.
        \item{\tt copies} -- number of copies to be produced from
        the selected objects.
        \end{itemize}

\item{\tt ci\_copyrotate2(bx, by, angle, copies, editaction )}
        \begin{itemize}
        \item{\tt bx, by} -- base point for rotation
        \item{\tt angle} -- angle by which the selected objects are incrementally
        shifted to make each copy.  {\tt angle} is measured in degrees.
        \item{\tt copies} -- number of copies to be produced from
        the selected objects.
        \item{\tt editaction}  0 --nodes, 1 -- lines (segments), 2 --block labels, 3 -- arc
                segments, 4- group
        \end{itemize}

\item{\tt ci\_copytranslate(dx, dy, copies)}
        \begin{itemize}
        \item{\tt dx,dy} -- distance by which the selected objects are incrementally shifted.
        \item{\tt copies} -- number of copies to be produced from the selected objects.
        \end{itemize}

\item{\tt ci\_copytranslate2(dx, dy, copies, editaction)}
        \begin{itemize}
        \item{\tt dx,dy} -- distance by which the selected objects are incrementally shifted.
        \item{\tt copies} -- number of copies to be produced from the selected objects.
        \item{\tt editaction}  0 --nodes, 1 -- lines (segments), 2 --block labels, 3 -- arc
                segments, 4- group
        \end{itemize}

\item{\tt ci\_createradius(x, y, r)} turns a corner located at {\tt (x,y)} into a curve of radius {\tt r}.

\item{\tt ci\_moverotate(bx,by,shiftangle)}
        \begin{itemize}
        \item{\tt bx, by} -- base point for rotation
        \item{\tt shiftangle} -- angle in degrees by which the selected objects are rotated.
        \end{itemize}

\item{\tt ci\_moverotate2(bx,by,shiftangle, editaction)}
        \begin{itemize}
        \item{\tt bx, by} -- base point for rotation
        \item{\tt shiftangle} -- angle in degrees by which the selected objects are rotated.
        \item{\tt editaction}  0 --nodes, 1 -- lines (segments), 2 --block labels, 3 -- arc
                segments, 4- group
        \end{itemize}

\item{\tt ci\_movetranslate(dx,dy)}
        \begin{itemize}
        \item{\tt dx,dy} -- distance by which the selected objects are shifted.
        \end{itemize}

\item{\tt ci\_movetranslate2(dx,dy,editaction)}
        \begin{itemize}
        \item{\tt dx,dy} -- distance by which the selected objects are shifted.
        \item{\tt editaction}  0 --nodes, 1 -- lines (segments), 2 --block labels, 3 -- arc
                segments, 4- group
        \end{itemize}

\item{\tt ci\_scale(bx,by,scalefactor)}
        \begin{itemize}
        \item{\tt bx, by} -- base point for scaling
        \item{\tt scalefactor} -- a multiplier that determines how
        much the selected objects are scaled
        \end{itemize}

\item{\tt ci\_scale2(bx,by,scalefactor,editaction)}
        \begin{itemize}
        \item{\tt bx, by} -- base point for scaling
        \item{\tt scalefactor} -- a multiplier that determines how
        much the selected objects are scaled
        \item{\tt editaction}  0 --nodes, 1 -- lines (segments), 2 --block labels, 3 -- arc
                segments, 4- group
        \end{itemize}

\item{\tt ci\_mirror(x1,y1,x2,y2)}
mirror the selected objects about a line passing through the points
{\tt (x1,y1)} and {\tt (x2,y2)}.

\item{\tt ci\_mirror2(x1,y1,x2,y2,editaction)}
mirror the selected objects about a line passing through the points
{\tt (x1,y1)} and {\tt (x2,y2)}. Valid {\tt editaction} entries are
0 for nodes, 1 for lines (segments), 2 for block labels, 3 for arc
segments, and 4 for groups.


\item \texttt{ci\_seteditmode(editmode)} Sets the current editmode to:

\texttt{"nodes"} -- nodes

\texttt{"segments"} - line segments

\texttt{"arcsegments"} - arc segments

\texttt{"blocks"} - block labels

\texttt{"group"} - selected group

This command will affect all subsequent uses of the other editing
commands, if they are used WITHOUT the \texttt{editaction}
parameter.
\end{itemize}

\subsection{Zoom Commands}

\begin{itemize}
\item \texttt{ci\_zoomnatural()} zooms to a ``natural'' view with sensible extents.

\item \texttt{ci\_zoomout()} zooms out by a factor of 50{\%}.

\item \texttt{ci\_zoomin()} zoom in by a factor of 200{\%}.

\item \texttt{ci\_zoom(x1,y1,x2,y2)} Set the display area to be from the bottom left
corner specified by \texttt{(x1,y1}) to the top right corner
specified by \texttt{(x2,y2)}.
\end{itemize}

\subsection{View Commands}
\begin{itemize}
\item \texttt{ci\_showgrid()} Show the grid points.

\item \texttt{ci\_hidegrid()} Hide the grid points points.

\item \texttt{ci\_gridsnap("flag")} Setting flag to ''on'' turns on snap to grid,
setting flag to "off"turns off snap to grid.

\item{\tt ci\_snapgridon} turns on snap-to-grid. Equivalent to {\tt ci\_gridsnap('on')}.

\item{\tt ci\_snapgridoff} turns off snap-to-grid. Equivalent to {\tt ci\_gridsnap('off')}.

\item \texttt{ci\_setgrid(density,"type")} Change the grid spacing. The density
parameter specifies the space between grid points, and the type
parameter is set to \texttt{"cart"} for Cartesian coordinates or
\texttt{"polar"} for polar coordinates.

\item \texttt{ci\_refreshview()} Redraws the current view.

\item{\tt ci\_minimize()} minimizes the active magnetics input view.

\item{\tt ci\_maximize()} maximizes the active magnetics input view.

\item{\tt ci\_restore()} restores the active magnetics input view from a
 minimized or maximized state.

\item{\tt ci\_resize(width,height)} resizes the active magnetics input
 window client area to width $\times$ height.
\item{\tt ci\_getview} grabs the currently displayed current flow input view and
 returns it as an image matrix, in the form \texttt{imread} would return.

\end{itemize}




\subsection{Object Properties}

\begin{itemize}
\item \texttt{ci\_getmaterial('materialname')} fetches the material specified by \texttt{materialname} 
from the materials library.

\item \texttt{ci\_addmaterial("materialname", ox, oy, ex, ey, ltx, lty)} adds a new material with
called \texttt{"materialname"} with the material properties:

\texttt{ox} Electrical conductivity in the x- or r-direction in units of S/m.

\texttt{oy} Electrical conductivity in the y- or z-direction in units of S/m.

\texttt{ex} Relative permittivity in the x- or r-direction.

\texttt{ey} Relative permittivity in the y- or z-direction.

\texttt{ltx} Dielectric loss tangent in the x- or r-direction.

\texttt{lty} Dielectric loss tangent in the y- or z-direction.

\item \texttt{ci\_addpointprop("pointpropname",Vp,jp)} adds a new point property of
name \texttt{"pointpropname"} with either a specified potential
\texttt{Vp }a point current density \texttt{jp} in units of A/m.

\item \texttt{ci\_addboundprop("boundpropname", Vs, js, c0, c1, BdryFormat)} adds a
new boundary property with name \texttt{"boundpropname"}

For a ``Fixed Voltage'' type boundary condition, set the
\texttt{Vs} parameter to the desired voltage and all other
parameters to zero.

To obtain a ``Mixed'' type boundary condition, set \texttt{C1} and
\texttt{C0} as required and \texttt{BdryFormat} to 1. Set all other
parameters to zero.

To obtain a prescribes surface current density, set \texttt{js} to
the desired current density in A/m$^{2}$ and set \texttt{BdryFormat}
to 2.

For a ``Periodic'' boundary condition, set \texttt{BdryFormat} to 3
and set all other parameters to zero.

For an ``Anti-Perodic'' boundary condition, set \texttt{BdryFormat}
to 4 set all other parameters to zero.

\item \texttt{ci\_addconductorprop("conductorname", Vc, jc, conductortype)} adds a new
conductor property with name \texttt{"conductorname"} with either a
prescribed voltage or a prescribed total current. Set the unused
property to zero. The \texttt{conductortype} parameter is 0 for
prescribed current and 1 for prescribed voltage.

\item \texttt{ci\_deletematerial("materialname")} deletes the material named
\texttt{"materialname"}.

\item \texttt{ci\_deleteboundprop("boundpropname")} deletes the boundary property
named \texttt{"boundpropname"}.

\item \texttt{ci\_deleteconductor("conductorname")} deletes the conductor named
\texttt{conductorname}.

\item \texttt{ci\_deletepointprop("pointpropname")} deletes the point property named
\texttt{"pointpropname"}

\item \texttt{ci\_modifymaterial("BlockName",propnum,value)} This function allows for
modification of a material's properties without redefining the
entire material (e.g. so that current can be modified from run to
run). The material to be modified is specified by
\texttt{"BlockName"}. The next parameter is the number of the
property to be set. The last number is the value to be applied to
the specified property. The various properties that can be modified
are listed below:

\begin{tabular}{lll}
propnum & Symbol &  Description \\ \hline
 0 & \texttt{BlockName} & Name of the material \\
 1 & \texttt{ox} & x (or r) direction conductivity \\
 2 & \texttt{oy} & y (or z) direction conductivity \\
 3 & \texttt{ex} & x (or r) direction relative permittivity \\
 4 & \texttt{ey} & y (or z) direction relative permittivity \\
 5 & \texttt{ltx} & x (or r) direction dielectric loss tangent \\
 6 & \texttt{lty} & y (or z) direction dielectric loss tangent \\
\end{tabular}


\item \texttt{ci\_modifyboundprop("BdryName",propnum,value) }This function allows for
modification of a boundary property. The BC to be modified is specified by
\texttt{"BdryName"}. The next parameter is the number of the property to be
set. The last number is the value to be applied to the specified property.
The various properties that can be modified are listed below:

\begin{tabular}{lll}
 \texttt{propnum} & Symbol & Description \\ \hline
 0 & \texttt{BdryName} & Name of boundary property \\
 1 & \texttt{Vs} & Fixed Voltage \\
 2 & \texttt{js} & Prescribed current density \\
 3 & \texttt{c0} & Mixed BC parameter \\
 4 & \texttt{c1} & Mixed BC parameter \\
 5 & \texttt{BdryFormat} & Type of boundary condition:\\
   &   &
   \begin{tabular}{lll}
   0 & = & Prescribed V \\
   1 & = & Mixed \\
   2 & = & Surface current density \\
   3 & = & Periodic \\
   4 & = & Antiperiodic
   \end{tabular}
\end{tabular}





\item \texttt{ci\_modifypointprop("PointName",propnum,value)} This function allows for
modification of a point property. The point property to be modified
is specified by \texttt{"PointName"}. The next parameter is the
number of the property to be set. The last number is the value to
be applied to the specified property. The various properties that
can be modified are listed below:



\begin{tabular}{lll}
\texttt{propnum} & Symbol & Description \\ \hline
 0 & \texttt{PointName} & Name of the point property \\
 1 & \texttt{Vp} & Prescribed nodal voltage \\
 2 & \texttt{jp} & Point current density in A/m
\end{tabular}

\item \texttt{ci\_modifyconductorprop("ConductorName",propnum,value)} This function
allows for modification of a conductor property. The conductor
property to be modified is specified by \texttt{"ConductorName"}.
The next parameter is the number of the property to be set. The
last number is the value to be applied to the specified property.
The various properties that can be modified are listed below:


\begin{tabular}{lll}
\texttt{propnum} & Symbol & Description \\ \hline
  0 & \texttt{ConductorName} & Name of the conductor property \\
  1 & \texttt{Vc} &  Conductor voltage \\
  2 & \texttt{jc} &  Total conductor current \\
  3 & \texttt{ConductorType} & 0 = Prescribed current, 1 = Prescribed voltage
\end{tabular}

\end{itemize}

\subsection{Miscellaneous}

\begin{itemize}
\item \texttt{ci\_savebitmap("filename")} saves a bitmapped screenshot of the current
view to the file specified by \texttt{"filename"}, subject to the
\texttt{printf}-type formatting explained previously for the
\texttt{ci\_saveas} command.

\item \texttt{ci\_savemetafile("filename")} saves a metafile screenshot of the current
view to the file specified by \texttt{"filename"}, subject to the
\texttt{printf}-type formatting explained previously for the
\texttt{ci\_saveas} command.

\item \texttt{ci\_refreshview()} Redraws the current view.

\item \texttt{ci\_close()} closes the preprocessor window and
destroys the current document.

\item \texttt{ci\_shownames(flag)} This function allow the user to display or hide the
block label names on screen. To hide the block label names,
\texttt{flag} should be 0. To display the names, the parameter
should be set to 1.

\item{\tt ci\_hidenames} hides the block label names. Equivalent to {\tt ci\_shownames(0)}.

\item{\tt ci\_gettitle} returns the file path of the current document as a string.

\item{\tt ci\_setcomment('comment')} sets the comment/description string associated with the current problem.

\item{\tt ci\_readdxf("filename")} This function imports a dxf file specified by {\tt "filename"}.

\item{\tt ci\_savedxf("filename")} This function saves geometry information in a dxf file specified by {\tt "filename"}.

\item{\tt ci\_defineouterspace(Zo,Ro,Ri)} defines
an axisymmetric external region to be used in conjuction with the
Kelvin Transformation method of modeling unbounded problems.  The
{\tt Zo} parameter is the z-location of the origin of the outer region,
the {\tt Ro} parameter is the radius of the outer region, and the {\tt
Ri} parameter is the radius of the inner region ({\em i.e.} the region of
interest). In the exterior region, the permeability varies as a function of
distance from the origin of the external region.  These parameters
are necessary to define the permeability variation in the external
region.

\item{\tt ci\_attachouterspace()} marks all
selected block labels as members of the external region used for
modeling unbounded axisymmetric problems via the Kelvin
Transformation.

\item{\tt ci\_detachouterspace()} undefines all selected block labels
as members of the external region used for modeling unbounded axisymmetric
problems via the Kelvin Transformation.

\item{\tt ci\_attachdefault()} marks the
selected block label as the default block label.  This block label
is applied to any region that has not been explicitly labeled.

\item{\tt ci\_detachdefault()} undefines the default
attribute for the selected block labels.

\item{\tt ci\_makeABC(n,R,x,y,bc)} creates a series of circular shells that emulate the
impedance of an unbounded domain (i.e. an Improvised Asymptotic Boundary
Condition).  The {\tt n} parameter contains the number of shells to be used
(should be between 1 and 10), {\tt R} is the radius of the solution domain, and
{\tt (x,y)} denotes the center of the solution domain.  The {\tt bc} parameter should
be specified as 0 for a Dirichlet outer edge or 1 for a Neumann outer edge.
If the function is called without all the parameters, the function makes up
reasonable values for the missing parameters. 

\end{itemize}

\section{Current Flow Post Processor Command Set}

There are a number of Lua scripting commands designed to operate in
the postprocessor. As with the preprocessor commands, these
commands can be used with either the underscore naming or with the
no-underscore naming convention.

\subsection{Data Extraction Commands}
\begin{itemize}

\item \texttt{co\_lineintegral(type)} Calculate the line integral for the defined contour

% inttype==0 => E.t 
% inttype==1 => J.n
% inttype==2 => Contour Length
% inttype==3 => Average Voltage
% inttype==4 => Stress Tensor Force
% inttype==5 => Stress Tensor Torque

\begin{tabular}{ll}
\texttt{type} & Integral \\ \hline
 0 & $E \cdot t$ \\
 1 & $J \cdot n$ \\
 2 & Contour length \\
 3 & Average voltage over contour \\
 4 & Force from stress tensor \\
 5 & Torque from stress tensor
\end{tabular}

This integral returns either 1 or 2 values, depending on the
integral type, $e.g.:$

\texttt{Fx, Fy = co\_lineintegral(4)}

\item \texttt{co\_blockintegral(type)} Calculate a block integral for the selected
blocks

\begin{tabular}{ll}
\texttt{type} & Integral \\ \hline
	 0  & Real Power \\
	 1  & Reactive Power \\
	 2  & Apparent Power \\
	 3  & Time-Average Stored Energy \\
	 4  & Block cross-section area \\
	 5  & Block volume \\
	 6  & x (or r) direction Weighted Stress Tensor force, DC component \\
	 7  & y (or z) direction Weighted Stress Tensor force, DC component \\
	 8  & x (or r) direction Weighted Stress Tensor force, 2x frequency component \\
	 9  & y (or z) direction Weighted Stress Tensor force, 2x frequency component \\
	 10 & Weighted Stress Tensor torque, DC component \\
	 11 & Weighted Stress Tensor torque, 2x frequency component
\end{tabular}

Returns a value that can be complex, as necessary.

\item \texttt{co\_getpointvalues(X,Y)} Get the values associated with the point at x,y
The return values, in order, are:

\begin{tabular}{ll}
\texttt{Symbol} &  Definition \\ \hline
\texttt{V} & Voltage \\
\texttt{Jx} & x- or r- direction component of current density \\
\texttt{Jy} & y- or z- direction component of current density \\
\texttt{Kx} & x- or r- direction component of AC conductivity \\
\texttt{Ky} & y- or z- direction component of AC conductivity \\
\texttt{Ex} & x- or r- direction component of electric field intensity \\
\texttt{Ey} & y- or z- direction component of electric field intensity \\
\texttt{ex} & x- or r- direction component of permittivity \\
\texttt{ey} & y- or z- direction component of permittivity \\
\texttt{Jdx} & x- or r- direction component of displacement current density \\
\texttt{Jdy} & y- or z- direction component of displacement current density \\
\texttt{ox} & x- or r- direction component of permittivity \\
\texttt{oy} & y- or z- direction component of permittivity \\
\texttt{Jcx} & x- or r- direction component of conduction current density \\
\texttt{Jcy} & y- or z- direction component of conduction current density \\
\end{tabular}

\item{\tt co\_getv(x,y)} returns the voltage at {\tt (x,y)} -- column 1 of {\tt co\_getpointvalues}. An equivalent form is {\tt co\_getv([x,y])}.

\item{\tt co\_getj(x,y)} returns the two components of current density {\tt Jx, Jy} at {\tt (x,y)} -- columns 2 and 3 of {\tt co\_getpointvalues}. An equivalent form is {\tt co\_getj([x,y])}.

\item{\tt co\_getk(x,y)} returns the two components of AC conductivity {\tt Kx, Ky} at {\tt (x,y)} -- columns 4 and 5 of {\tt co\_getpointvalues}. An equivalent form is {\tt co\_getk([x,y])}.

\item{\tt co\_gete(x,y)} returns the two components of electric field intensity {\tt Ex, Ey} at {\tt (x,y)} -- columns 6 and 7 of {\tt co\_getpointvalues}. An equivalent form is {\tt co\_gete([x,y])}.

\item \texttt{co\_makeplot(PlotType,NumPoints,Filename,FileFormat) }Allows Lua access
to the X-Y plot routines. If only PlotType or only PlotType and NumPoints
are specified, the command is interpreted as a request to plot the requested
plot type to the screen. If, in addition, the Filename parameter is
specified, the plot is instead written to disk to the specified file name as
an extended metafile. If the FileFormat parameter is also, the command is
instead interpreted as a command to write the data to disk to the specfied
file name, rather than display it to make a graphical plot. Valid entries
for PlotType are:

\begin{tabular}{ll}
\texttt{PlotType} &  Definition \\ \hline
 0 & \texttt{V (Voltage)} \\
 1 & \texttt{$\vert$J$\vert$ (Magnitude of current density)} \\
 2 & \texttt{J . n (Normal current density)} \\
 3 & \texttt{J . t (Tangential current density)} \\
 4 & \texttt{$\vert$E$\vert$ (Magnitude of field intensity)} \\
 5 & \texttt{E . n (Normal field intensity)} \\
 6 & \texttt{E . t (Tangential field intensity)} \\
 7 & \texttt{$\vert$Jc$\vert$ (Magnitude of conduction current density)} \\
 8 & \texttt{Jc . n (Normal conduction current density)} \\
 9 & \texttt{Jc . t (Tangential conduction current density)} \\
 10 & \texttt{$\vert$Jd$\vert$ (Magnitude of displacement current density)} \\
 11 & \texttt{Jd . n (Normal displacement current density)} \\
 12 & \texttt{Jd . t (Tangential displacement current density)} \\
\end{tabular}

\vspace*{8pt}
Valid file formats are:

\begin{tabular}{ll}
\texttt{FileFormat} &  Definition  \\ \hline
 0 & \texttt{Multi-column text with legend } \\
 1 & \texttt{Multi-column text with no legend } \\
 2 & \texttt{Mathematica-style formatting}
\end{tabular}

For example, if one wanted to plot $V $ to the screen with 200 points evaluated
to make the graph, the command would be: \newline
\texttt{co\_makeplot(0,200) } \newline
If this plot were to be written to disk as a metafile, the command
would be:\texttt{ } \newline
\texttt{co\_makeplot(0,200,"c:temp.emf")} \newline
To write data instead of a plot to disk, the command would be of
the form:\texttt{ } \newline
\texttt{co\_makeplot(0,200,"c:temp.txt",0)} \newline


\item \texttt{co\_getprobleminfo() } 
Returns info on problem description. Returns four values:

\begin{tabular}{ll} \hline
Return value & Definition \\ \hline 
1 &  problem type  \\ 
2 &  frequency in Hz \\ 
3 &  depth assumed for planar problems in meters \\
4 &  length unit used to draw the problem, represented in meters.
\end{tabular}

\item \texttt{co\_getconductorproperties("conductor")} Properties are returned for the
conductor property named "conductor". Two values are returned: The voltage
of the specified conductor, and the current on the specified
conductor.
\end{itemize}

\subsection{Selection Commands}

\begin{itemize}
\item \texttt{co\_seteditmode(mode) }Sets the mode of the postprocessor to point,
contour, or area mode. Valid entries for mode are\texttt{ "point",
"contour", }and\texttt{ "area". }

\item \texttt{co\_selectblock(x,y) }Select the block that contains point\texttt{
(x,y).}

\item \texttt{co\_groupselectblock(n) }Selects all of the blocks that are labeled by
block labels that are members of group n. If no number is specified
({\em i.e.} \texttt{co\_groupselectblock()}), all blocks are selected.

\item \texttt{co\_selectconductor("name")} Selects all nodes, segments, and arc
segments that are part of the conductor specified by the string
\texttt{("name")}. This command is used to select conductors for the
purposes of the ``weighted stress tensor'' force and torque
integrals, where the conductors are points or surfaces, rather than
regions ({\em i.e.} can't be selected with \texttt{co\_selectblock}).

\item \texttt{co\_addcontour(x,y) }Adds a contour point at\texttt{ (x,y). }If this is
the first point then it starts a contour, if there are existing
points the contour runs from the previous point to this point. The
\texttt{co\_addcontour }command has the same functionality as a
right-button-click contour point addition when the program is
running in interactive mode.\texttt{ }

\item \texttt{co\_bendcontour(angle,anglestep)} Replaces the straight line formed by
the last two points in the contour by an arc that spans angle
degrees. The arc is actually composed of many straight lines, each
of which is constrained to span no more than anglestep degrees. The
\texttt{angle} parameter can take on values from -180 to 180
degrees. The \texttt{anglestep} parameter must be greater than
zero. If there are less than two points defined in the contour,
this command is ignored.

\item \texttt{co\_selectpoint(x,y)} Adds a contour point at the closest input point to
\texttt{(x,y).} If the selected point and a previous selected points lie at
the ends of an arcsegment, a contour is added that traces along the
arcsegment. The \texttt{selectpoint} command has the same
functionality as the left-button-click contour point selection when
the program is running in interactive mode.

\item \texttt{co\_clearcontour()} Clear a prevously defined contour

\item \texttt{co\_clearblock()} Clear block selection
\end{itemize}



\subsection{Zoom Commands}

\begin{itemize}
\item \texttt{co\_zoomnatural()} Zoom to the natural boundaries of the geometry.

\item \texttt{co\_zoomin()} Zoom in one level.

\item \texttt{co\_zoomout()} Zoom out one level.

\item \texttt{co\_zoom(x1,y1,x2,y2)} Zoom to the window defined by lower left corner
\texttt{(x1,y1)} and upper right corner\texttt{ (x2,y2).}
\end{itemize}



\subsection{View Commands}

\begin{itemize}
\item \texttt{co\_showmesh()} Show the mesh.

\item \texttt{co\_hidemesh()} Hide the mesh.

\item \texttt{co\_showpoints()} Show the node points from the input geometry.

\item \texttt{co\_hidepoints()} Hide the node points from the input geometry.

\item \texttt{co\_smooth("flag")} This function controls whether or not smoothing is
applied to the $D$ and $E$ fields which are naturally piece-wise
constant over each element. Setting flag equal to \texttt{"on"}
turns on smoothing, and setting flag to \texttt{"off"} turns off
smoothing.

\item{\tt co\_smoothon} turns on smoothing of the $J$ and $E$ fields, which are naturally piece-wise constant over each element. Equivalent to {\tt co\_smooth('on')}.

\item{\tt co\_smoothoff} turns off smoothing. Equivalent to {\tt co\_smooth('off')}.

\item \texttt{co\_showgrid()} Show the grid points.

\item \texttt{co\_hidegrid()} Hide the grid points points.

\texttt{co\_gridsnap("flag")} Setting flag to ''on'' turns on snap to grid,
setting flag to "off" turns off snap to grid.

\item \texttt{co\_setgrid(density,"type")} Change the grid spacing. The density
parameter specifies the space between grid points, and the type
parameter is set to \texttt{"cart"} for Cartesian coordinates or
\texttt{"polar"} for polar coordinates.

\item \texttt{co\_hidedensityplot()} hides the current density plot.

\item \texttt{co\_showdensityplot(legend,gscale,type,upper,lower)} Shows the
current density plot with options:

\texttt{legend} Set to 0 to hide the plot legend or 1 to show the plot
legend.

\texttt{gscale} Set to 0 for a colour density plot or 1 for a grey scale
density plot.

\texttt{upper} Sets the upper display limit for the density plot.

\texttt{lower} Sets the lower display limit for the density plot.

\texttt{type} Sets the type of density plot. Specific choices for the type of density plot include:

\begin{tabular}{ll}
\texttt{type} &  Description  \\ \hline
0 & $|V|$		\\
1 & $| \mbox{Re}(V)|$	\\
2 & $| \mbox{Im}(V)|$	\\
3 & $|J|$		\\
4 & $|\mbox{Re}(J)|$	\\
5 & $|\mbox{Im}(J)|$	\\
6 & $|E|$		\\
7 & $|\mbox{Re}(E)|$	\\
8 & $|\mbox{Im}(E)|$
\end{tabular}

\item \texttt{co\_hidecontourplot()} Hides the contour plot.

\item \texttt{co\_showcontourplot(numcontours,lower{\_}V,upper{\_}V),type} shows the
$V$ contour plot with options:

\texttt{numcontours} Number of equipotential lines to be plotted;

\texttt{upper{\_}V} Upper limit for contours;

\texttt{lower{\_}V} Lower limit for contours; 

\texttt{type} the type of contour plot to be rendered.

If \texttt{co\_numcontours} is -1 all parameters are ignored and
default values are used, \\ e.g. \texttt{show{\_}contour{\_}plot(-1)}

The type can take on the values of {\tt "real"}, {\tt "imag"}, or {\tt "both"}, denoting
the real part of voltage, the imaginary part of voltage, or both components of voltage.

\item \texttt{co\_showvectorplot(type,scalefactor)}
controls the display of vectors denoting the field strength and
direction. The \texttt{type} parameter can take on the following values:

\begin{tabular}{ll}
\texttt{type} &  Description  \\ \hline
0 & No vector plot \\
1 & $\mbox{Re}(J)$ \\
2 & $\mbox{Re}(E)$ \\
3 & $\mbox{Im}(J)$ \\
4 & $\mbox{Im}(E)$ \\
5 & $\mbox{Re}(J)$ and $Im(J)$ \\
6 & $\mbox{Re}(E)$ and $Im(E)$
\end{tabular}
The \texttt{scalefactor}
determines the relative length of the vectors. If the scale is set
to 1, the length of the vectors are chosen so that the highest field
magnitude corresponds to a vector that is the same length as the
current grid size setting.

\item{\tt co\_minimize()} minimizes the active magnetics input view.

\item{\tt co\_maximize()} maximizes the active magnetics input view.

\item{\tt co\_restore()} restores the active magnetics input view from a
 minimized or maximized state.

\item{\tt co\_resize(width,height)} resizes the active magnetics input
 window client area to width $\times$ height.
\item{\tt co\_getview} grabs the currently displayed current flow output view and
 returns it as an image matrix, in the form \texttt{imread} would return.

\end{itemize}

\subsection{Miscellaneous}

\begin{itemize}

\item \texttt{co\_close()} close the current postprocessor window.

\item{\tt co\_refreshview()} Redraws the current view.

\item{\tt co\_reload()} Reloads the solution from disk.

\item \texttt{co\_savebitmap("filename")} saves a bitmapped screen shot of the current
view to the file specified by \texttt{"filename"}. Note that if you
use a path you must use two backslashes (e.g.
\texttt{"c:$\backslash \backslash $temp$\backslash \backslash
$myfile.bmp"}). If the file name contains a space (e.g. file
names like c:$\backslash $program files$\backslash $stuff) you must
enclose the file name in (extra) quotes by using a
\texttt{$\backslash $"} sequence. For example:

\texttt{co\_savebitmap("$\backslash $"c:$\backslash \backslash
$temp$\backslash \backslash $screenshot.bmp$\backslash $"")}

\item \texttt{co\_savemetafile("filename")} saves a metafile screenshot of the current
view to the file specified by \texttt{"filename"}, subject to the
printf-type formatting explained previously for the
\texttt{savebitmap} command.

\item \texttt{co\_shownames(flag)} This function allow the user to display or hide the
block label names on screen. To hide the block label names,
\texttt{flag} should be 0. To display the names, the parameter
should be set to 1.

\item{\tt co\_hidenames} hides the block label names. Equivalent to {\tt co\_shownames(0)}.

\item{\tt co\_gettitle} returns the file path of the current document as a string.

\item{\tt co\_snapgrid('flag')} setting {\tt flag} to {\tt 'on'} turns on snap to grid, setting it to {\tt 'off'} turns it off.

\item{\tt co\_snapgridon} turns on snap-to-grid. Equivalent to {\tt co\_snapgrid('on')}.

\item{\tt co\_snapgridoff} turns off snap-to-grid. Equivalent to {\tt co\_snapgrid('off')}.

\item{\tt co\_numnodes} Returns the number of nodes in the in focus current flow output mesh.
\item{\tt co\_numelements} Returns the number of elements in the in focus current flow output mesh.
\item{\tt co\_getnode(n)} Returns the (x,y) or (r,z) position of the nth mesh node.
\item{\tt co\_getelement(n)} returns the following properties for the nth element:
    \begin{enumerate}
        \item Index of first element node
        \item Index of second element node
        \item Index of third element node
        \item x (or r) coordinate of the element centroid
        \item y (or z) coordinate of the element centroid
        \item element area using the length unit defined for the problem
        \item group number associated with the element
    \end{enumerate}


\item {\bf Bulk getters.}  These fetch the whole mesh in one call instead of one item at a time, which is much faster than looping the per-item getters.  Mesh node indices are {\bf 1-based}, matching {\tt co\_getnode} and {\tt co\_getelement}.  Each returns:

\begin{itemize}
\item{\tt co\_getallnodes} an $N \times 2$ matrix of the coordinates of every mesh node.
\item{\tt co\_getallelements} an $N \times 7$ matrix with one row per mesh element, as {\tt co\_getelement}.
\item{\tt co\_getallj} an $N \times 2$ matrix of the current density $J_x, J_y$ at every element. Complex at a nonzero frequency. The value is constant over each element; no smoothing is applied and no point search is performed.
\item{\tt co\_getallv} the voltage at every mesh node, complex at a nonzero frequency. Note that this is per {\em node}, whereas {\tt co\_getallj} is per {\em element}.
\end{itemize}

\end{itemize}


%-----------------------------------------------------------------------------------

\end{document}
